% Chapter 2. All quoted numbers come from chapters/ch02-numbers.tex, written by code/ch02/run_all.py.
% Generated by code/ch02/run_all.py. Do not edit by hand.
\newcommand{\altFeeCost}{8.5}
\newcommand{\altFeeFlat}{20.0}
\newcommand{\altFeeStep}{12.55}
\newcommand{\altFlatTotal}{105}
\newcommand{\annIncOne}{16.0}
\newcommand{\annMgmtShare}{55.6}
\newcommand{\annNetGrowth}{6.4}
\newcommand{\annPV}{283.0}
\newcommand{\annPVi}{125.8}
\newcommand{\annPVm}{157.2}
\newcommand{\annPVperp}{1,000}
\newcommand{\annPVtwenty}{486.0}
\newcommand{\annPhi}{3.6}
\newcommand{\bClawCap}{8.40}
\newcommand{\bClawback}{6.15}
\newcommand{\bDbdCarry}{14.00}
\newcommand{\bDealACarry}{9.00}
\newcommand{\bDealAHurdle}{37.79}
\newcommand{\bDealAProfit}{45}
\newcommand{\bDealBCarry}{3.00}
\newcommand{\bDealBHurdle}{34.01}
\newcommand{\bDealBProfit}{15}
\newcommand{\bDealCCarry}{0.00}
\newcommand{\bDealCHurdle}{29.39}
\newcommand{\bDealCProfit}{\ensuremath{-}20}
\newcommand{\bDealDCarry}{2.00}
\newcommand{\bDealDHurdle}{15.87}
\newcommand{\bDealDProfit}{10}
\newcommand{\bFeeIP}{10.00}
\newcommand{\bFeePost}{0.75}
\newcommand{\bFeeTotal}{10.75}
\newcommand{\bFeeYearFive}{0.45}
\newcommand{\bFeeYearSeven}{0.15}
\newcommand{\bFeeYearSix}{0.15}
\newcommand{\bGrossIrr}{10.99}
\newcommand{\bInvested}{85}
\newcommand{\bLPirrDbdClaw}{8.45}
\newcommand{\bLPirrDbdNoClaw}{7.24}
\newcommand{\bLPirrWf}{9.55}
\newcommand{\bPaidIn}{95.75}
\newcommand{\bPoolA}{9.00}
\newcommand{\bPoolB}{3.00}
\newcommand{\bPoolC}{0.00}
\newcommand{\bPoolCarry}{12.00}
\newcommand{\bPoolClaw}{4.15}
\newcommand{\bPoolD}{0.00}
\newcommand{\bPoolEntFour}{10.00}
\newcommand{\bPoolEntOne}{9.00}
\newcommand{\bPoolEntThree}{8.00}
\newcommand{\bPoolEntTwo}{12.00}
\newcommand{\bProceeds}{135}
\newcommand{\bProfit}{39.25}
\newcommand{\bTVPIgross}{1.41}
\newcommand{\bTVPInet}{1.33}
\newcommand{\bWfCarry}{7.85}
\newcommand{\bWfCatchA}{0.00}
\newcommand{\bWfCatchB}{0.00}
\newcommand{\bWfCatchC}{0.00}
\newcommand{\bWfCatchD}{6.01}
\newcommand{\bWfGPA}{0.00}
\newcommand{\bWfGPB}{0.00}
\newcommand{\bWfGPC}{0.00}
\newcommand{\bWfGPD}{7.85}
\newcommand{\bWfLPA}{75.00}
\newcommand{\bWfLPB}{40.00}
\newcommand{\bWfLPC}{0.00}
\newcommand{\bWfLPD}{12.15}
\newcommand{\bWfPrefA}{0.00}
\newcommand{\bWfPrefB}{19.55}
\newcommand{\bWfPrefC}{0.00}
\newcommand{\bWfPrefD}{4.50}
\newcommand{\bWfRocA}{75.00}
\newcommand{\bWfRocB}{20.45}
\newcommand{\bWfRocC}{0.00}
\newcommand{\bWfRocD}{0.30}
\newcommand{\bWfShare}{20.0}
\newcommand{\bWfSplitA}{0.00}
\newcommand{\bWfSplitB}{0.00}
\newcommand{\bWfSplitC}{0.00}
\newcommand{\bWfSplitD}{9.19}
\newcommand{\carryFullFifteen}{2.39}
\newcommand{\carryFullFortyFive}{8.06}
\newcommand{\carryFullThirty}{5.50}
\newcommand{\carryFullThirtyRich}{7.65}
\newcommand{\carryHardFifteen}{1.07}
\newcommand{\carryHardFortyFive}{6.29}
\newcommand{\carryHardThirty}{3.65}
\newcommand{\carryMCThirty}{5.49}
\newcommand{\carryMCn}{1,000,000}
\newcommand{\carryMCse}{0.012}
\newcommand{\carryMCz}{\ensuremath{-}0.37}
\newcommand{\carryNoHurdleThirty}{6.40}
\newcommand{\carryVegaHigh}{\ensuremath{-}0.059}
\newcommand{\carryVegaHighThirty}{0.114}
\newcommand{\carryVegaLow}{0.240}
\newcommand{\carryVstar}{158.67}
\newcommand{\exFourCloseOne}{114.80}
\newcommand{\exFourCloseThree}{105.92}
\newcommand{\exFourCloseTwo}{95.86}
\newcommand{\exFourIncOne}{3.70}
\newcommand{\exFourIncThree}{0.00}
\newcommand{\exFourIncTwo}{0.00}
\newcommand{\exOneA}{29.0}
\newcommand{\exOneB}{171}
\newcommand{\exSevenRatio}{1.0513}
\newcommand{\exSevenVegaDown}{0.067}
\newcommand{\exSevenVegaPeak}{0.077}
\newcommand{\exSixPV}{54.73}
\newcommand{\exThreeGP}{10.00}
\newcommand{\exThreeGPHard}{2.79}
\newcommand{\exThreeH}{136.05}
\newcommand{\exThreeVstar}{145.06}
\newcommand{\fofSingleOne}{37.3}
\newcommand{\fofSingleZero}{32.2}
\newcommand{\fofTotalOne}{55.5}
\newcommand{\fofTotalZero}{52.8}
\newcommand{\fofUnderOne}{37.3}
\newcommand{\fofUnderZero}{37.4}
\newcommand{\hfCloseFour}{119.29}
\newcommand{\hfCloseOne}{110.40}
\newcommand{\hfCloseThree}{102.98}
\newcommand{\hfCloseTwo}{97.15}
\newcommand{\hfGrossCagr}{7.62}
\newcommand{\hfGrossEnd}{134.14}
\newcommand{\hfGrossFour}{123.58}
\newcommand{\hfGrossOne}{115.00}
\newcommand{\hfGrossThree}{104.92}
\newcommand{\hfGrossTwo}{99.36}
\newcommand{\hfHardIncFour}{1.13}
\newcommand{\hfHardIncOne}{1.60}
\newcommand{\hfHwmFour}{119.29}
\newcommand{\hfHwmOne}{110.40}
\newcommand{\hfHwmThree}{110.40}
\newcommand{\hfHwmTwo}{110.40}
\newcommand{\hfIncFour}{2.22}
\newcommand{\hfIncOne}{2.60}
\newcommand{\hfIncThree}{0.00}
\newcommand{\hfIncTwo}{0.00}
\newcommand{\hfMgmtFour}{2.06}
\newcommand{\hfMgmtOne}{2.00}
\newcommand{\hfMgmtThree}{1.94}
\newcommand{\hfMgmtTwo}{2.21}
\newcommand{\hfNetCagr}{4.51}
\newcommand{\hfNetFour}{15.84}
\newcommand{\hfNetOne}{10.40}
\newcommand{\hfNetThree}{6.00}
\newcommand{\hfNetTwo}{\ensuremath{-}12.00}
\newcommand{\hfNoHwmIncThree}{1.17}
\newcommand{\hfNoHwmTotal}{15.62}
\newcommand{\hfOpenFour}{102.98}
\newcommand{\hfOpenOne}{100.00}
\newcommand{\hfOpenThree}{97.15}
\newcommand{\hfOpenTwo}{110.40}
\newcommand{\hfPreFour}{121.52}
\newcommand{\hfPreOne}{113.00}
\newcommand{\hfPreThree}{102.98}
\newcommand{\hfPreTwo}{97.15}
\newcommand{\hfSoftIncFour}{2.22}
\newcommand{\hfSoftIncOne}{2.60}
\newcommand{\hfTotalFees}{13.03}
\newcommand{\mcCapBindN}{0}
\newcommand{\mcClawProb}{99.9}
\newcommand{\mcDbdCarry}{17.85}
\newcommand{\mcDbdIrrCost}{0.53}
\newcommand{\mcDist}{160.5}
\newcommand{\mcFeesMean}{13.15}
\newcommand{\mcGrossIrr}{11.05}
\newcommand{\mcGrossNetSd}{10.9}
\newcommand{\mcMeanMult}{1.92}
\newcommand{\mcN}{6,000}
\newcommand{\mcNetIrr}{9.58}
\newcommand{\mcNetTVPI}{1.48}
\newcommand{\mcOverCommitN}{0}
\newcommand{\mcOwed}{6.49}
\newcommand{\mcOwedGivenPos}{6.49}
\newcommand{\mcOwedNinetyFive}{13.11}
\newcommand{\mcPaidIn}{97.2}
\newcommand{\mcPaidMax}{99.67}
\newcommand{\mcPoolCarry}{15.60}
\newcommand{\mcPoolClawProb}{93.2}
\newcommand{\mcPoolFeeCluster}{37.0}
\newcommand{\mcPoolIrrCost}{0.47}
\newcommand{\mcPoolOwed}{4.24}
\newcommand{\mcPoolOwedNinetyFive}{10.50}
\newcommand{\mcPreFeeIrr}{14.71}
\newcommand{\mcTVPI}{1.59}
\newcommand{\mcWedgeCarry}{1.52}
\newcommand{\mcWedgeFees}{3.65}
\newcommand{\mcWfAboveDbd}{0.0}
\newcommand{\mcWfCarry}{11.36}
\newcommand{\mcWfOverAll}{0.13}
\newcommand{\mcWfOverMean}{0.35}
\newcommand{\mcWfOverNinetyFive}{0.33}
\newcommand{\mcWfOverProb}{36.9}
\newcommand{\mcWfPaid}{11.49}
\newcommand{\mcZeroDbd}{0.1}
\newcommand{\mcZeroPool}{6.8}
\newcommand{\mcZeroWf}{34.7}
\newcommand{\modModCloseFive}{112.24}
\newcommand{\modModCloseFour}{102.02}
\newcommand{\modModCloseSix}{123.01}
\newcommand{\modModCloseThree}{92.07}
\newcommand{\modModCloseTwo}{83.10}
\newcommand{\modModIncFive}{2.03}
\newcommand{\modModIncFour}{1.10}
\newcommand{\modModIncSix}{2.69}
\newcommand{\modModIncThree}{1.00}
\newcommand{\modModIncTwo}{0.90}
\newcommand{\modModSum}{7.72}
\newcommand{\modTarget}{106.25}
\newcommand{\modTradCloseFive}{114.31}
\newcommand{\modTradCloseFour}{104.30}
\newcommand{\modTradCloseSix}{125.28}
\newcommand{\modTradCloseThree}{94.08}
\newcommand{\modTradCloseTwo}{84.00}
\newcommand{\modTradIncFive}{2.50}
\newcommand{\modTradIncFour}{1.07}
\newcommand{\modTradIncSix}{2.74}
\newcommand{\modTradIncThree}{0.00}
\newcommand{\modTradIncTwo}{0.00}
\newcommand{\modTradSum}{6.32}
\newcommand{\netEffRate}{33.5}
\newcommand{\netFeeA}{7.20}
\newcommand{\netFeeB}{1.50}
\newcommand{\netGross}{10.00}
\newcommand{\netPaid}{4.35}
\newcommand{\netSingle}{3.20}
\newcommand{\optAtSigTen}{1.12}
\newcommand{\optAtSigTwenty}{1.88}
\newcommand{\optDownSigForty}{1.42}
\newcommand{\optDownSigTwenty}{0.31}
\newcommand{\optExactAtSigTwenty}{1.92}
\newcommand{\optUpSigTwenty}{6.73}
\newcommand{\optVegaAt}{0.077}
\newcommand{\optVegaDown}{0.044}
\newcommand{\optVegaUp}{0.031}
\newcommand{\progCommitForHundred}{15.15}
\newcommand{\progExposure}{9.35}
\newcommand{\progNavSteady}{6.60}
\newcommand{\progNetCash}{0.79}
\newcommand{\progOvercommit}{1.42}
\newcommand{\progUncalledSteady}{2.75}
\newcommand{\reBalOne}{10.75}
\newcommand{\reBalTwo}{2.39}
\newcommand{\reDist}{17.86}
\newcommand{\reEarlyInv}{2.04}
\newcommand{\reInvIrr}{12.32}
\newcommand{\reInvTot}{15.39}
\newcommand{\reMulInvIrr}{11.13}
\newcommand{\reMulSpoTot}{3.23}
\newcommand{\reP}{15}
\newcommand{\reProfit}{7.86}
\newcommand{\reProjIrr}{13.36}
\newcommand{\rePromote}{0.69}
\newcommand{\reSpoIrr}{21.25}
\newcommand{\reSpoShareProfit}{18.7}
\newcommand{\reSpoTot}{2.47}
\newcommand{\reTenIrr}{7.27}
\newcommand{\reTenSpoIrrTiers}{1.79}
\newcommand{\reTenSpoMulTiers}{3.23}
\newcommand{\reTierOneInv}{12.79}
\newcommand{\reTierOneSize}{11.94}
\newcommand{\reTierOneSpo}{1.42}
\newcommand{\reTierThreeInv}{0.21}
\newcommand{\reTierThreeSize}{0.33}
\newcommand{\reTierThreeSpo}{0.12}
\newcommand{\reTierTwoInv}{2.39}
\newcommand{\reTierTwoSize}{3.32}
\newcommand{\reTierTwoSpo}{0.93}
\newcommand{\reYrFive}{15.60}
\newcommand{\shareDet}{48.4}
\newcommand{\shareHwmForty}{67.5}
\newcommand{\shareHwmTen}{49.8}
\newcommand{\shareHwmTwenty}{54.8}
\newcommand{\shareN}{40,000}
\newcommand{\shareNoHwmForty}{80.0}
\newcommand{\shareNoHwmTen}{53.1}
\newcommand{\shareNoHwmTwenty}{62.5}
\newcommand{\spCatch}{11.73}
\newcommand{\spGP}{16.00}
\newcommand{\spGPHard}{6.61}
\newcommand{\spGPNoHurdle}{16.00}
\newcommand{\spGPPart}{16.00}
\newcommand{\spGrossIrr}{12.47}
\newcommand{\spH}{146.93}
\newcommand{\spLP}{164.00}
\newcommand{\spLPirr}{10.40}
\newcommand{\spPref}{46.93}
\newcommand{\spRoc}{100.00}
\newcommand{\spSplitGP}{4.27}
\newcommand{\spSplitLP}{17.07}
\newcommand{\spSplitTotal}{21.33}
\newcommand{\spV}{180}
\newcommand{\spVstarFull}{158.67}
\newcommand{\spVstarPart}{162.58}
\newcommand{\taBreakeven}{8}
\newcommand{\taCalledFive}{0.94}
\newcommand{\taCalledTotal}{1.000}
\newcommand{\taDistTotal}{1.79}
\newcommand{\taPeakBowHigh}{1.30}
\newcommand{\taPeakBowLow}{0.80}
\newcommand{\taPeakBowMid}{1.04}
\newcommand{\taPeakNav}{1.04}
\newcommand{\taPeakYear}{5}
\newcommand{\taTVPI}{1.79}
\newcommand{\taTVPIGHigh}{2.64}
\newcommand{\taTVPIGLow}{1.21}
\newcommand{\taTVPIGMid}{1.79}
\newcommand{\taTroughCash}{0.79}
\newcommand{\taTroughYear}{4}

\begin{chapterrefs}
\setcounter{chapter}{1}
\chapter{Fund Structures, Fees and Distribution Waterfalls}
\label{ch:structures}

An investor in a private fund does not own the assets. It owns an interest in a vehicle, and the
contract that creates the vehicle decides who controls the assets, when cash moves, and how gains are
shared between the investors and the manager. Two funds that hold identical portfolios can deliver very
different returns to their investors because their contracts differ. This chapter studies those
contracts with the same care that Chapter~1 gave to return measurement: it describes the legal
vehicles and the terms that matter, then turns the main fee and profit sharing rules into formulas, proves
the properties that shape the manager's incentives, and checks each result with code.

The return an investor receives is what is left after the money has passed through the contract,
from the first capital call to the final distribution. The chapter follows it along that route, for
one fund and then for \mcN{} simulated ones, and then asks what the fees are worth to the manager and
what they encourage it to do. The same tools then serve for property promotes and for planning
commitments.

\begin{keybox}[Reproducing the numbers]
The scripts in \texttt{code/ch02/} write the numbers they compute into a file of macros
that the chapter reads, so no result is copied by hand; contract parameters such as an 8\% hurdle are
stated assumptions, and the few figures taken from published sources are cited where they appear. Tests check the scripts against closed forms, Monte
Carlo estimates, an independent IRR library and cross-checks run in R and GNU Octave. The simulations use NumPy's
\texttt{default\_rng} with seeds 202 to 205, one for each simulation, and each seed is
given where its simulation is described.
\end{keybox}
% =====================================================================================
\section{Why the contract matters as much as the portfolio}
\label{sec2:why}

A fund is a contract between people who supply capital and people who supply judgement. The investors\ix{Investor}
cannot watch every decision, and the manager knows more than they do about its own skill and, once it is
investing, about the deals and its effort. After the money is
committed the manager may serve itself, by taking too much risk, charging expenses to the fund or
spending its time on a successor fund\ix{Successor fund}; economists call this hidden action \emph{moral hazard}\index{Moral hazard}. Before
it is committed, investors cannot tell good managers from bad ones, and the managers keenest for
capital on easy terms may be the weakest; this hidden information creates \emph{adverse selection}\index{Adverse selection}. Fund contracts respond in a few ways. Carried
interest\ix{Carried interest} and incentive fees tie pay to performance, and a manager that invests its own money shares
in the losses. Investors also get rights to information, to vote on conflicts and, in the last
resort, to remove the manager.

The same tools create new problems. A share of profits with no share of losses is an option, and
the holder of a plain call gains from volatility; \cref{sec2:options} shows when carry does not. A fee charged on assets rewards asset growth whether or not the
assets earn a return. Much of this chapter is about measuring those side effects precisely, because a
carry term that costs investors little when returns are steady can cost them a lot when returns are
volatile. Chapter~1 showed how to measure what a
fund earned; this chapter shows how much of it reaches the investor.

\FloatBarrier
% =====================================================================================
\section{What does an investor actually own?}
\label{sec2:vehicles}

Before a fund buys anything, its sponsor has to choose a legal wrapper, and the choice is a compromise.
Investors\ix{Investor} want to lose no more than they commit, they want the fund's income taxed once and in their
own hands, and they want terms they can bargain over. No single vehicle gives every investor all of
this, which is why one fund often uses several.

The \emph{limited partnership}\index{Limited partnership} is the usual answer for private equity, private credit and real estate
funds, and for hedge funds\ix{Hedge fund} sold to domestic investors. Its \emph{limited partners\ix{Limited partner}}\index{Limited partner} (LPs) put up most of
the capital and can lose that capital and any commitment still outstanding, but in general no more,
subject to guarantees they give and to local law. One \emph{general partner\ix{General partner}}\index{General partner} (GP) runs the fund and
carries unlimited liability. Sponsors blunt that exposure by making the GP a limited liability\ix{Limited liability} company, so claims against the GP
are ordinarily limited to the company's assets. That does not shield the principals from guarantees
they sign, from their own misconduct or from statutory exceptions\ix{Exceedance}, and the details depend on the
governing law.

Limited partners generally leave the running of the fund to the GP, and how far they may go without
risking their protection depends on local law. In many places an LP that takes part in managing the
fund can lose its limited liability, and each jurisdiction's statute decides where voting on a
conflict ends and taking part in management begins. In return the partnership is usually looked
through for tax\ix{Taxes}, so gains are taxed once in the partners' hands, and its terms come from an agreement
the parties write for themselves. How much either is worth depends on the jurisdiction and on the investor.

A company changes the bargain. Shareholders generally have limited liability with no duty to stay passive\ix{Index fund},
subject to statutory exceptions, personal guarantees and their own misconduct,
but in many jurisdictions the company is taxed as an entity, with the outcome depending on its
classification\ix{Classification} and on the investor. For a taxable domestic investor that can mean a second layer of tax, depending on the local rules for
taxing dividends\ix{Dividends} and gains.
For a pension plan\ix{Pension plan} or a foreign investor it can be an advantage, because income passed through a
partnership may count as business income that such an investor would prefer not to receive directly.
Offshore hedge funds are therefore usually companies (\cref{fig2:structure}). The \emph{limited
liability company} (LLC) sits between the two, with limited liability for every member, flexible
internal rules and, in the United States where the form began, pass-through tax by default; other
countries may tax it as a company. Canada, for one, generally treats a US LLC as a corporation for its
own tax purposes, with consequences that depend on the investor and on the tax treaty. Sponsors use it for the general partner,
for the management company and to hold single assets.

A manager with both kinds of investor can still run one portfolio. In the \emph{master-feeder}\index{Master-feeder} structure
of \cref{fig2:structure}, each group enters through the feeder that suits its tax position, an onshore
partnership or an offshore company, and both feeders put all their capital into a master fund. Trades
happen once, at the master, and are allocated pro rata, so the feeders earn the same gross return;
their net returns still differ by the costs, taxes and fees charged in each feeder.

Smaller vehicles do narrower jobs. A \emph{special purpose vehicle}\index{Special purpose vehicle} (SPV) holds one property, one loan
pool or one co-investment, so that the asset and its debts sit apart from the rest of the sponsor's
business. How far apart raises two separate questions: would the asset stay out of the sponsor's
estate if the sponsor became insolvent, and must the sponsor consolidate the vehicle in its accounts?
Documents that confine the vehicle to its purpose, limit its debts to those tied to its assets and
require an independent director's vote for a bankruptcy filing are commonly used to make it
\emph{bankruptcy remote}\index{Bankruptcy remote}, though how well they work depends on ownership, guarantees and the
insolvency law that applies. They lower the risk that the sponsor's insolvency reaches the asset
without removing it. The securitisations of Chapter~6 depend on the first question: a true sale, and
a vehicle that a court will not merge with the sponsor in its insolvency. Whether the sponsor
consolidates the vehicle in its accounts is a separate accounting matter, and many consolidated
vehicles serve their investors well.

\begin{figure}[t]
  \centering
  \includegraphics[width=0.92\textwidth]{ch02/fig01_structure}
  \caption{A master-feeder fund. Two feeders with different tax treatment invest in one master fund.
  The general partner controls the partnership and receives the carried interest\index{Carried interest} or incentive
  allocation; the management company advises the fund and receives the management fee\index{Management fee}.}
  \label{fig2:structure}
\end{figure}

The split between the general partner and the management company in \cref{fig2:structure} is
deliberate. The management fee is income for services and is paid to the management company. The
carried interest\ix{Carried interest} is a share of partnership profits and is allocated to the general partner. Its tax
treatment depends on the fund's jurisdiction, the applicable rules and how the allocation is
structured. Service providers sit around the structure: a
prime broker or custodian holds the assets, an administrator keeps the books and computes the net asset
value, and an auditor reports on the financial statements.

\FloatBarrier
% =====================================================================================
\section{How do investors get in and out?}
\label{sec2:openclosed}

\paragraph{Open-end funds}
An \emph{open-end}\index{Open-end} or \emph{evergreen}\index{Evergreen} fund has no fixed life. Investors\ix{Investor} subscribe and redeem at the net
asset value\ix{Net asset value} (NAV) on stated dealing dates, usually monthly or quarterly, after a notice period\ix{Notice period}. Because
investors can leave, the fund must manage liquidity\ix{Liquidity risk} against its redemption\ix{Redemption} terms, with liquid
assets, cash reserves, borrowing or contractual limits on redemptions. Strategies in
liquid markets can do that, which is why most hedge funds\ix{Hedge fund} are open-end. Liquidity terms protect the remaining investors from the cost of
forced sales\ix{Liquidation}: a \emph{lock-up\ix{Lock-up}}\index{Lock-up} bars redemptions for an initial period (a soft lock-up allows them for a
fee), a \emph{gate} caps the share of the fund that can be redeemed on one date, and a \emph{side pocket}\index{Side pocket}
holds illiquid positions apart. An investor who redeems usually keeps its side pocket interest and is
paid when those positions are sold, on terms the documents set. The documents also usually let the
fund suspend redemptions in a crisis.

In an open-end fund investors arrive at different times and at different NAVs. The
\emph{high-water mark\ix{High-water mark}}\index{High-water mark} means the manager earns an incentive\ix{Incentive compatibility} fee\ix{Incentive fee}\index{Incentive fee} only on gains that take the NAV above
its highest level at an earlier fee date, or above its starting level if that is higher. With one mark for the whole fund,
an investor who subscribes after a loss gets a free ride: it pays no incentive fee on its own gains
until the fund regains the old mark. If the fund falls from 100 to 80 and
an investor buys in at 80, a rise to 95 gives that investor a gain of almost 19\% and no incentive fee,
because the fund is still below its mark of 100. Funds remove this with series accounting, which keeps separate
accounting series for subscriptions with distinct fee histories, or with equalisation adjustments;
\cref{sec2:hf} develops the economics of the high-water mark.

\paragraph{Drawdown funds}
A \emph{drawdown\ix{Drawdown}} fund, also called a closed-end\ix{Drawdown fund} or private equity\ix{Private equity} style fund, raises
\emph{commitments\ix{Commitment}}: promises to provide capital up to a stated amount when asked. The general partner\ix{General partner}
\emph{calls} capital as it finds investments, usually during an \emph{investment period}\index{Investment period} of about five
years, then manages and sells the investments during a \emph{harvest period}\index{Harvest period}, and distributes the
proceeds. The fund has a fixed term, commonly ten years, with options to extend for one or two years to
complete sales. Investors cannot redeem. They can sell their interest to another investor on the
secondary market\ix{Secondary market}, usually at a negotiated price and with the general partner's consent. An LP that
fails to meet a capital call is a \emph{defaulting} investor and faces penalties set in the partnership
agreement, which may include the forfeiture of part of its interest.

\paragraph{Risks particular to drawdown funds}
A drawdown fund exposes its investors to several risks that are larger than in a liquid public
market investment. The investor commits
to a \emph{blind pool}\index{Blind pool}, because the investments are not known when the commitment is made. It must keep
liquidity available for calls of uncertain size and timing, which creates \emph{funding risk}\index{Funding risk}. It has no
contractual redemption right, although a secondary sale may offer an exit at a negotiated price, so
it bears most of the illiquidity of the assets. Its returns depend on one team, so the departure
of key people or a change of strategy is a large risk. And reported values in the early years are
estimates that can be slow to reflect bad news, which, as Chapter~1 showed, can smooth the return series.
Investors answer these risks chiefly by spreading commitments across managers and vintage years\ix{Vintage year} and by
planning them with a cash flow model (\cref{sec2:pacing}); due diligence\ix{Due diligence}, key person\ix{Key person clause} clauses and
secondary sales add some protection, though not complete protection.

\paragraph{How fees differ}
Hedge fund and private equity fees share a common outline, a recurring management fee\ix{Management fee} and a fee or allocation linked to performance,
but differ in the details that matter. A hedge fund usually charges its management fee on NAV and its
incentive fee on gains above a high-water mark, usually crystallised and paid each year, though some
funds crystallise more often and some less. A private equity
fund usually charges its management fee on committed capital during the investment period and on
invested capital afterwards, and earns carried interest\ix{Carried interest} only when gains are realised, usually by sales, on cumulative
profits, usually after a preferred return. The private equity design measures performance over the whole life
of the fund, and that is one reason for its more complex waterfall.

\FloatBarrier
% =====================================================================================
\section{What does the management fee cost over a fund's life?}
\label{sec2:mgmtfees}

Let $K$ be committed capital\ix{Commitment}, $f_1$ the fee rate during an investment period of $T_I$ years and $f_2$
the rate afterwards, charged on the cost of investments not yet realised, $I_t$. With fees charged at
the start of each year $t = 0, 1, \dots, T-1$,
\begin{equation}\label{eq2:mgmtfee}
  F_t = \begin{cases} f_1 K, & t < T_I,\\ f_2\, I_t, & t \ge T_I,\end{cases}
  \qquad
  F = \sum_{t=0}^{T-1} F_t = f_1 K T_I + f_2 \sum_{t=T_I}^{T-1} I_t .
\end{equation}
The fee is paid by calling capital from the LPs\ix{Limited partner}, so fees use up commitment. With no recycling\ix{Recycling} or borrowing, at most $K - F$ of the commitment can be invested. Under the stated schedule the fee on committed capital is fixed for the
investment period once the commitment is made, barring contractual reductions, suspension or early
termination, and it is charged even on capital never invested. A fee on invested capital stops
charging for capital once it is returned, but because it falls as investments are sold it gives the
manager a reason to hold assets longer.

\begin{example}{Fund B: management fees}\label{ex2:fees}
Exits occur at the start of the year shown, before that year's fee is charged. Fund B has
commitments of 100 and makes the four investments in \cref{tab2:deals}, at a total cost of
\bInvested. It charges 2\% of commitments for the first five years and 1.5\% of the cost of
unrealised deals afterwards. The fees are $5\times 2 = \bFeeIP$ for years 0 to 4. At the start of
year 5 deals C and D are unrealised, with cost 30, so the fee is $0.015\times 30 = \bFeeYearFive$. In
years 6 and 7 only deal D remains and the fee is \bFeeYearSix{} each year; the fund ends with the sale
of D in year 8. Total fees are \bFeeTotal. The LPs pay in \bPaidIn, the cost of the deals plus the
fees, and the rest of the commitment is never called; the fund distributes \bProceeds. Panel (a) of
\cref{fig2:fees} shows the fee base
dropping from commitments to invested cost\ix{Cost basis}.
\end{example}

\begin{table}[t]
  \centering
  \caption{The four investments of Fund B. Proceeds are received at the start of the year of exit,
  before that year's fee.}
  \label{tab2:deals}
  \begin{tabular}{lrrrrr}
    \toprule
    Deal & Year invested & Cost & Year of exit & Proceeds & Profit \\
    \midrule
    A & 0 & 30 & 3 & 75 & \bDealAProfit \\
    B & 1 & 25 & 5 & 40 & \bDealBProfit \\
    C & 1 & 20 & 6 & 0  & \bDealCProfit \\
    D & 2 & 10 & 8 & 20 & \bDealDProfit \\
    \bottomrule
  \end{tabular}
\end{table}

Fee terms that look similar can differ a lot over a fund's life. Panel (b) of \cref{fig2:fees} compares
three sets of terms for a fund that invests 17 a year for five years and holds each investment for five
years, so the cost of unrealised deals is 17, 34, 51, 68 and 85 in years 0 to 4 and then falls by 17 a
year. A flat 2\% on commitments for ten years costs \altFeeFlat{} per 100 committed; stepping down to
1.5\% of cost after the investment period costs \altFeeStep; charging 2\% of invested cost throughout
costs \altFeeCost. Under the flat terms fees and investments together reach \altFlatTotal, more than the
commitment, so the comparison treats the excess as paid out of proceeds, which an agreement can allow;
it compares fee terms and leaves capital budgets aside. The difference between the first and the last, more than ten percent of commitments,
is charged whatever the investments earn. \textcite{phalippou2009} documents how fees of this kind, together
with transaction and monitoring fees\ix{Transaction fee} charged to portfolio companies, can materially reduce gross
private equity returns in his sample of funds.

\begin{figure}[t]
  \centering
  \includegraphics[width=\textwidth]{ch02/fig02_fees}
  \caption{(a) Fund B: the fee base is commitments during the investment period and the cost of
  unrealised deals afterwards; bars show the annual fee multiplied by ten. (b) Cumulative fees per 100 of
  commitments under three fee terms for a fund that invests evenly over five years.}
  \label{fig2:fees}
\end{figure}

\FloatBarrier
% =====================================================================================
\section{How are distributions split?}
\label{sec2:waterfall}

\subsection{When have the investors earned the hurdle?}

The \emph{distribution waterfall\ix{Distribution waterfall}}\index{Distribution waterfall} is the clause of the partnership agreement that splits each
distribution between the LPs\ix{Limited partner} and the GP\ix{General partner}. Proceeds flow through a sequence of tiers, and each tier fills
before the next receives anything. The waterfall used in this chapter, a common
private equity design, has four.
\begin{enumerate}[label=(\roman*)]
  \item \emph{Return of capital}: everything goes to the LPs until they have received back all capital
        contributed, including capital called to pay fees and expenses.
  \item \emph{Preferred return\ix{Preferred return}}: everything goes to the LPs until they have also received a return at the
        \emph{hurdle rate\ix{Hurdle rate}}\index{Hurdle rate} $h$ on that capital, compounded; the examples use 8\% a year, a common
        choice, though each agreement sets its own rate and some funds have no preferred return.
  \item \emph{Catch-up\ix{Catch-up}}: a share $k$ of each dollar, often all of it, goes to the GP until the GP has
        received its carried interest share $c$, commonly 20\%, of all profits distributed so far.
  \item \emph{Split}: the remainder is shared $1-c$ to the LPs and $c$ to the GP.
\end{enumerate}
The first two tiers can be stated together with one quantity.

\begin{definition}[Hurdle balance]\label{def2:balance}
Let the LPs contribute $C_i$ at times $t_i$ and receive distributions $L_j$ at times $s_j$, over a
finite history, and let $h>-1$. The \emph{hurdle balance}\index{Hurdle balance} at time $t$ is
\begin{equation}\label{eq2:balance}
  B(t) = \sum_{i:\,t_i\le t} C_i\,(1+h)^{t-t_i} \;-\; \sum_{j:\,s_j< t} L_j\,(1+h)^{t-s_j},
\end{equation}
the amount the LPs would need to receive at $t$ to have earned exactly $h$ on their flows to date.
The balance is computed immediately before the distribution at $t$ is allocated: contributions made
at $t$ count, and the distribution being allocated at $t$ does not. The balance can be
negative. It models the LPs' preferred return as the waterfall above defines it; agreements word it in different ways.
\end{definition}

\begin{proposition}\label{prop2:balance}
$B(t) \le 0$ if and only if the LPs' flows counted in \cref{eq2:balance}, that is, contributions up to
and including $t$ and distributions strictly before $t$, have a nonnegative net present value\ix{Net present value} at rate $h$.
If those flows are conventional (at least one contribution, all contributions before all
distributions, and at least one positive distribution), this holds if and only if the LPs' internal
rate of return on those same flows is at least $h$.
\end{proposition}
\begin{proof}
Multiplying \cref{eq2:balance} by $(1+h)^{-t} > 0$ gives
$(1+h)^{-t}B(t) = \sum_i C_i (1+h)^{-t_i} - \sum_j L_j (1+h)^{-s_j} = -\NPV_h$, the negative of the NPV
at rate $h$ of the LPs' flows with contributions counted as outflows. So $B(t)\le 0 \iff \NPV_h \ge 0$.
For conventional flows, with the same cutoff, let $\tau$ be the date of the last contribution, which by
assumption precedes every distribution counted. Then $f(r) = (1+r)^{\tau}\NPV_r = -\sum_i C_i(1+r)^{\tau-t_i} + \sum_j L_j(1+r)^{\tau-s_j}$
has the same sign as $\NPV_r$, and, given at least one positive distribution, it is strictly
decreasing on $r > -1$ because every exponent
$\tau - t_i$ is nonnegative and every exponent $\tau - s_j$ is negative. As $r\downarrow-1$ the distribution terms grow without bound, so $f\to+\infty$, and as $r\to\infty$ they
vanish while the contribution terms stay at or below $-C_\tau<0$, so $f$ is eventually negative. Hence
$f$ has exactly one root $r^*$, the IRR, as in Chapter~1, and $\NPV_h \ge 0 \iff f(h) \ge 0 \iff h \le r^*$.
\end{proof}

Tiers (i) and (ii) therefore pay the LPs $\min\{X,\max\{B(t),0\}\}$ out of a distribution $X$. The split
of that amount between capital and preferred return matters only for reporting: it is capital first,
up to the unreturned contributions, then preferred return. The balance counts only what the LPs receive; carry paid to the GP does not enter it. Once the LPs have been paid in the split tier, $B(t)$ is negative,
and in this sequential allocation that cushion absorbs later fee calls: a fee called after the hurdle
sends the next distribution back to tiers (i) and (ii) only if it exceeds the cushion. The cushion
decides only whether the LPs regain priority. Any call after the catch-up is complete raises the
paid-in capital\ix{Paid-in capital} $P$ of \cref{prop2:catchup} without changing $D$ or the GP's carry $G$, so cumulative
profit $D-P$ falls and the GP at once holds more than $c$ of it. Split tier dollars keep that excess
(\cref{prop2:catchup}), and only the end-of-life true-up measured in \cref{sec2:mc} removes it. This is how the models in this chapter apply the
tiers; agreements that track unreturned capital separately can sequence distributions differently.

\subsection{How large is the catch-up?}

\begin{proposition}[Size of the catch-up]\label{prop2:catchup}
Suppose a distribution arrives after the hurdle\ix{Hurdle rate} is met. Let $P$ be the LPs\ix{Limited partner}' total contributions
(paid-in capital\ix{Paid-in capital} rather than commitments\ix{Commitment}), $D$ cumulative distributions to the LPs and the GP\ix{General partner} together,
including the part of this one paid in tiers (i) and (ii), and $G$ the GP's cumulative
carry\ix{Cumulative carry} so far. With catch-up\ix{Catch-up} rate $k > c$, the catch-up tier has size
\begin{equation}\label{eq2:catchup}
  y^{*} = \frac{c\,(D-P) - G}{k-c},
\end{equation}
when this is positive, and after it the GP's cumulative carry equals $c$ times cumulative profit,
$G + k y^{*} = c\,(D + y^{*} - P)$. Every later dollar paid in the split tier keeps that equality.
\end{proposition}
\begin{proof}
In the catch-up tier a dollar raises the GP's carry by $k$ and cumulative profit by one. The tier ends
when $G + k y = c\,(D + y - P)$, which is linear in $y$ with solution \cref{eq2:catchup}; $k > c$ makes
the denominator positive, so the GP's share of profit rises towards $c$ as $y$ grows. In the split tier
a dollar raises carry by $c$ and profit by one, so if $G' = c\,(D'-P)$ before it, then
$G' + c x = c\,(D' + x - P)$ after it.
\end{proof}

With a full catch-up ($k = 1$), one contribution $K$ and one distribution, so that $P=K$ and
$D = H$, \cref{eq2:catchup} becomes
$y^{*} = c(H-K)/(1-c)$, where $H$ is the hurdle amount. Each dollar of preferred return paid to the LPs
must be matched by $c/(1-c)$ dollars of catch-up, or 0.25 with $c = 20\%$, to restore an 80/20
split of profit.

\subsection{Soft and hard hurdles}

In the structures modelled here, the catch-up\ix{Catch-up} decides what kind of hurdle\ix{Hurdle rate} the fund has. With a catch-up the hurdle is \emph{soft}: it
delays the carry but, once the fund is far enough above it, the GP receives $c$ of all profit, including
the profit below the hurdle. Without a catch-up the hurdle is \emph{hard}: the GP earns carry only on
profit above the hurdle amount. The difference is easiest to see in the simplest fund.

\begin{proposition}[One contribution, one distribution]\label{prop2:single}
Let the LPs contribute $K$ at time 0 and the fund distribute $V$ at time $T$, so the hurdle amount is
$H = K(1+h)^T$. With catch-up rate $k>c$, define $V^{*} = (kH - cK)/(k-c)$. Then the carried interest is
\begin{equation}\label{eq2:gpcases}
  G(V) = \begin{cases}
    0, & V \le H,\\[2pt]
    k\,(V-H), & H < V \le V^{*},\\[2pt]
    c\,(V-K), & V > V^{*},
  \end{cases}
\end{equation}
and with a hard hurdle $G(V) = c\,(V-H)^{+}$, where $x^{+} = \max\{x,0\}$.
\end{proposition}
\begin{proof}
For $V\le H$ tiers (i) and (ii) absorb everything. Above $H$ the catch-up pays the GP $k$ per dollar.
By \cref{prop2:catchup} with $P=K$, $D = H$ and $G = 0$ it ends after $y^* = c(H-K)/(k-c)$, that is at
\begin{equation*}
  V = H + \frac{c(H-K)}{k-c} = \frac{(k-c)H + cH - cK}{k-c} = \frac{kH - cK}{k-c} = V^{*}.
\end{equation*} At $V^{*}$ the GP
holds $k y^{*} = c(V^{*}-K)$, and each further dollar adds $c$, so $G(V) = c(V-K)$ for $V > V^{*}$. With a
hard hurdle tier (iii) is absent and tier (iv) starts at $H$.
\end{proof}

\begin{example}{A single-period fund}\label{ex2:single}
Take $K = 100$, $T = 5$, $h = 8\%$ and $c = 20\%$, so $H = 100 \times 1.08^5 = \spH$ and, with full
catch-up, $V^* = (\spH - 20)/0.8 = \spVstarFull$. Suppose the fund distributes gross proceeds
$V = \spV$, all of it available to the waterfall. Tier (i)
pays the LPs \spRoc, tier (ii) pays them \spPref, the catch-up pays the GP $0.25\times\spPref =
\spCatch$, and the remaining $\spV - \spVstarFull = \spSplitTotal$ is split \spSplitLP{} to the LPs and \spSplitGP{}
to the GP. The GP receives $\spGP = 0.2\times(\spV - 100)$, exactly 20\% of profit, and the LPs receive
\spLP. An 80\% catch-up gives the same \spGPPart, because $V$ is above that design's threshold of
\spVstarPart{} too. A hard hurdle gives only $0.2\times(\spV-\spH) = \spGPHard$. The fund's gross return
is \spGrossIrr\% a year and the LPs' net return is \spLPirr\% a year.
\end{example}

\Cref{fig2:tiers} shows the whole picture. Panel (a) stacks the four tiers against $V$; panel (b) plots
the GP's payoff under four designs. The kinks at $H$ and $V^*$ are where the incentives in
\cref{sec2:behaviour} come from. \Cref{lst2:core} is a compact implementation of the single-period case
that reproduces \cref{ex2:single}, including the hard hurdle case; the general version in \texttt{ch02lib.py} works with any dated flows.

\begin{figure}[t]
  \centering
  \includegraphics[width=\textwidth]{ch02/fig03_tiers}
  \caption{The single-period waterfall with $K = 100$, $T = 5$, $h = 8\%$ and $c = 20\%$. (a) How
  proceeds $V$ are paid through the tiers with a full catch-up. (b) Carried interest under a full
  catch-up, an 80\% catch-up, a hard hurdle and no hurdle.}
  \label{fig2:tiers}
\end{figure}

\Needspace*{21\baselineskip}
\lstinputlisting[language=Python,caption={The single-period waterfall (\texttt{code/ch02/snippets/waterfall\_core.py}).},label=lst2:core]{code/ch02/snippets/waterfall_core.py}

\FloatBarrier
% =====================================================================================
\section{What is the carry worth, and how does risk change it?}
\label{sec2:options}

\Cref{eq2:gpcases} is piecewise linear with slopes $0$, $k$ and $c$, which means it can be written with
call option payoffs.

\begin{proposition}[Option decomposition]\label{prop2:options}
Under the conditions of \cref{prop2:single},
\begin{equation}\label{eq2:decomp}
  G(V) = k\,(V-H)^{+} \;-\; (k-c)\,(V-V^{*})^{+},
\end{equation}
so the carried interest is $k$ calls struck at $H$ less $k-c$ calls struck at $V^{*}$. A hard hurdle is
$c$ calls struck at $H$, and no hurdle is $c$ calls struck at $K$.
\end{proposition}
\begin{proof}
Check the three regions of \cref{eq2:gpcases}. For $V \le H$ both terms are zero. For $H < V \le V^*$ the
second term is zero and the first is $k(V-H)$. For $V > V^*$ the right side is
$k(V-H) - (k-c)(V-V^*) = cV - kH + (k-c)V^*$, and $(k-c)V^* = kH - cK$, so it equals $c(V-K)$.
\end{proof}

If the value of the fund's portfolio $V_t$ followed a geometric Brownian motion with constant
volatility, paid nothing out before $T$ and could be traded continuously, and the interest rate were
constant, \cref{prop2:options} would give the value of the carry today by the Black-Scholes formula
derived in Chapter~5. Write $C(S,X,T,r,\sigma)$ for the price of a call on $S$ struck at $X$ and expiring at $T$.
Then
\begin{equation}\label{eq2:carryvalue}
  \mathcal{G}_0 = k\,C(V_0, H, T, r, \sigma) - (k-c)\,C(V_0, V^{*}, T, r, \sigma),
  \qquad C = S\,\Phi(d_1) - X e^{-rT}\Phi(d_2),
\end{equation}
with $d_{1,2} = [\ln(S/X) + (r \pm \sigma^2/2)T]/(\sigma\sqrt T)$. Private assets are not traded, so the replicating hedge behind the formula cannot be
built and \cref{eq2:carryvalue} gives only a model value. It still answers this section's question: how
the GP's claim responds to the risk of the portfolio.
\textcite{metrick2010} value the carry of a large sample of funds with a related option model.

\begin{example}{The value of carry}\label{ex2:carryvalue}
Take the fund of \cref{ex2:single} with $V_0 = 100$, so the portfolio is worth its cost, a risk-free
rate of 3\% and $T = 5$. With full catch-up the carry is worth \carryFullFifteen{} at 15\% volatility,
\carryFullThirty{} at 30\% and \carryFullFortyFive{} at 45\%. With a hard hurdle the values are
\carryHardFifteen, \carryHardThirty{} and \carryHardFortyFive; with no hurdle at 30\% it is
\carryNoHurdleThirty. A Monte Carlo estimate (seed 202) with \carryMCn{} antithetic draws gives \carryMCThirty{} at
30\% with a standard error of \carryMCse, within one standard error of the formula. If the GP's work
lifts the portfolio to $V_0 = 115$ the carry at 30\% is worth \carryFullThirtyRich. These are model values. Carry rarely trades, and when it does, as in
GP stakes deals, it is priced privately on other assumptions.
\end{example}

Panel (a) of \cref{fig2:carryoption} shows the value rising with volatility in each design, for a fund
whose portfolio is worth its cost. The sign is not guaranteed in general: the payoff is a difference of calls, so its
sensitivity to volatility can be negative at some portfolio values. Differentiating
\cref{eq2:carryvalue},
\begin{equation}\label{eq2:vega}
  \frac{\partial \mathcal{G}_0}{\partial \sigma}
  = \sqrt{T}\,V_0 \bigl[\, k\,\varphi\bigl(d_1(H)\bigr) - (k-c)\,\varphi\bigl(d_1(V^{*})\bigr) \bigr],
\end{equation}
where $\varphi$ is the standard normal density, because the vega of a call is $S\varphi(d_1)\sqrt T$.
When $V_0$ is far above $V^*$, $d_1(H) > d_1(V^*) > 0$, so $\varphi(d_1(H))$ is the smaller density and
the bracket can be negative: the short call struck at $V^*$ makes the payoff concave there. At 15\%
volatility the carry gains \carryVegaLow{} per point of volatility when $V_0 = 100$ but changes by
\carryVegaHigh{} per point when $V_0 = 200$, so there it loses value as volatility rises; at 30\% volatility and $V_0 = 200$ it gains
\carryVegaHighThirty. Far above the catch-up threshold the GP gains much less from extra risk, and at low volatility loses
from it, so it may prefer to protect its position; near or below the hurdle the gain from risk is
largest.
\textcite{carpenter2000} makes the same point for option pay in general: convex pay does not always
raise the appetite for risk.

\begin{figure}[t]
  \centering
  \includegraphics[width=\textwidth]{ch02/fig04_carry_option}
  \caption{Carried interest valued as options, with $K = 100$, $T = 5$, $h = 8\%$, $c = 20\%$ and a 3\%
  risk-free rate. (a) Value against volatility when the portfolio is worth its cost, with Monte Carlo
  checks. (b) Value against the current portfolio value; the dotted line is the value with zero
  volatility.}
  \label{fig2:carryoption}
\end{figure}

\FloatBarrier
% =====================================================================================
\section{Does it matter whether carry is measured deal by deal?}
\label{sec2:dbd}

\paragraph{Two ways to measure profit}
A \emph{whole-fund}\index{Distribution waterfall!European} or \emph{European} waterfall applies the tiers of \cref{sec2:waterfall} to the cash
flows of the fund as a whole: the LPs get back all contributions, for every deal and for fees, plus the
preferred return, before the GP receives any carry. A \emph{deal-by-deal}\index{Distribution waterfall!American} or \emph{American} waterfall
applies the tiers to each investment separately, so the GP is paid carry on a profitable exit even if
other investments are losing money. Many deal-by-deal agreements soften this by requiring that the
costs of deals already realised, including written-off deals, be recovered first. Investor guidance
generally favours the whole-fund form \parencite{ilpa2019}. Both descriptions are simplified
archetypes; agreements differ in how invested capital, fees and realised losses enter each test.

\begin{proposition}[Deal by deal pays at least as much without a hurdle]\label{prop2:subadd}
Let deals $i = 1,\dots,n$ have profits $\pi_i$ measured at a common date and no hurdle. Then
\begin{equation}\label{eq2:subadd}
  \sum_{i=1}^{n} c\,\pi_i^{+} \;\ge\; c\Bigl(\sum_{i=1}^{n}\pi_i\Bigr)^{+},
\end{equation}
with equality if and only if the nonzero profits all have the same sign.
\end{proposition}
\begin{proof}
The function $x\mapsto x^{+}$ is subadditive: $(a+b)^{+} \le a^{+} + b^{+}$, because $a+b \le a^{+}+b^{+}$
and $0 \le a^{+}+b^{+}$. Induction on $n$ gives $(\sum\pi_i)^{+} \le \sum\pi_i^{+}$; multiply by $c>0$. If
all $\pi_i \ge 0$ or all $\pi_i \le 0$ both sides agree. Otherwise some $\pi_i<0$, so $\sum\pi_i^{+} > \sum_i \pi_i$,
and $\sum\pi_i^{+} > 0$, so $\sum\pi_i^{+} > \max\{\sum\pi_i, 0\} = (\sum\pi_i)^{+}$.
\end{proof}

With hurdles and dated flows the comparison is no longer exact: the whole-fund hurdle compounds on
every contribution, fees and losing deals included, while a deal's hurdle compounds only on that
deal's cost. A contractual \emph{clawback}\index{Clawback} seeks to reconcile the overpayment at the end of the fund,
subject to any repayment caps and collection risk.

\begin{definition}[Clawback]\label{def2:clawback}
At the end of the fund, the GP must return to the LPs the excess of the carry it has received over the
carry it would have received under the whole-fund waterfall applied to the fund's complete cash flows,
\begin{equation}\label{eq2:clawback}
  \mathrm{CB} = \Bigl(\sum_{\text{exits}} G^{\mathrm{deal}} - G^{\mathrm{fund}}\Bigr)^{+},
\end{equation}
often capped at the carry received less the taxes the GP's principals paid on it, $(1-\tau)\sum G^{\mathrm{deal}}$, where $\tau$ is the principals' tax rate on carry.
This is the gross economic shortfall. The amount legally repayable can differ: agreements adjust for
taxes, cap the repayment, hold part of the carry in escrow and set their own allocation rules.
\end{definition}

A clawback is only as good as the GP's ability to pay it. Investors protect it with an \emph{escrow} or
holdback of part of each carry payment, with personal guarantees from the principals, and with interim
clawback tests during the fund's life rather than only at the end.

\begin{example}{Fund B under both waterfalls}\label{ex2:fundb}
Return to Fund B (\cref{ex2:fees,tab2:deals}). Its LP flows are the deal costs called in years 0, 1 and 2, the fees called at the start
of each year from 0 to 7 and the exit proceeds of \cref{tab2:deals}; they have an IRR of \bGrossIrr\% before
carry and a total profit of \bProfit. \Cref{tab2:fundb} applies the whole-fund waterfall. The sale of A in
year 3 and of B in year 5 only return capital and pay the preferred return; the hurdle balance is
cleared only by the sale of D in year 8, which pays the GP a catch-up of \bWfCatchD{} and leaves \bWfSplitD{} for the
split. The GP's carry is \bWfCarry, which is \bWfShare\% of the profit.

Under a pure deal-by-deal waterfall each deal is tested on its own cost and an 8\% hurdle on that cost.
Deal A cost 30 in year 0, so its hurdle amount is $30\times1.08^3 = \bDealAHurdle$. With a full catch-up the
GP has 20\% of the deal's whole profit once proceeds pass $(\bDealAHurdle - 0.2\times30)/0.8$, which 75 does,
so deal A pays carry of $0.2\times\bDealAProfit = \bDealACarry$ in year 3. That is carry paid early; the
clawback below is the GP's eventual liability. Deal B pays \bDealBCarry{} in year 5, C pays nothing and D, whose hurdle amount
is \bDealDHurdle, pays \bDealDCarry{} in year 8. The GP receives \bDbdCarry{} in total and owes a clawback
of $\bDbdCarry - \bWfCarry = \bClawback$, well within a cap net of tax, which at a 40\% tax rate is $0.6\times\bDbdCarry =
\bClawCap$. If realised losses are netted first, the GP's entitlement on the pool of realised deals is
\bPoolEntOne{} after A, \bPoolEntTwo{} after B, \bPoolEntThree{} after the write-off of C and \bPoolEntFour{}
after D. Carry already paid is never reduced before the end, so D pays nothing, the GP receives
\bPoolCarry{} and it owes \bPoolClaw.

The timing matters to the LPs even after the clawback has been paid. Their net IRR is \bLPirrWf\% under
the whole-fund waterfall, \bLPirrDbdClaw\% under the deal-by-deal waterfall with the clawback paid in
year 8, and \bLPirrDbdNoClaw\% if the clawback is never paid. The GP held \bDealACarry{} from year 3 and another \bDealBCarry{} from year 5, and received a further
\bDealDCarry{} in year 8, against \bWfCarry{} that it would have received only in year 8 under the
whole-fund waterfall; it returns the \bClawback{} excess in year 8 without interest.
\end{example}

\begin{table}[t]
  \centering
  \caption{Fund B under the whole-fund waterfall ($h = 8\%$, $c = 20\%$, full catch-up). Each row is one
  distribution, received at the start of the year shown; tier columns show the incremental amounts
  paid from that distribution, rounded to two decimals.}
  \label{tab2:fundb}
  \begin{tabular}{lrrrrrrr}
    \toprule
    Exit & Amount & Capital & Pref.\ return & Catch-up & Split & To LPs & To GP \\
    \midrule
    A, year 3 & 75 & \bWfRocA & \bWfPrefA & \bWfCatchA & \bWfSplitA & \bWfLPA & \bWfGPA \\
    B, year 5 & 40 & \bWfRocB & \bWfPrefB & \bWfCatchB & \bWfSplitB & \bWfLPB & \bWfGPB \\
    C, year 6 & 0  & \bWfRocC & \bWfPrefC & \bWfCatchC & \bWfSplitC & \bWfLPC & \bWfGPC \\
    D, year 8 & 20 & \bWfRocD & \bWfPrefD & \bWfCatchD & \bWfSplitD & \bWfLPD & \bWfGPD \\
    \bottomrule
  \end{tabular}
\end{table}

\begin{figure}[t]
  \centering
  \includegraphics[width=\textwidth]{ch02/fig05_fundb}
  \caption{Fund B. (a) Cumulative carry paid to the GP under the whole-fund waterfall and under
  deal-by-deal waterfalls with and without netting of realised losses; the vertical drop in year 8 is
  the clawback. (b) Contributions and distributions under the whole-fund waterfall.}
  \label{fig2:fundb}
\end{figure}

\FloatBarrier
% =====================================================================================
\section{What happens across many funds?}
\label{sec2:mc}

One fund proves little, so this section simulates \mcN{} funds. Each fund has commitments of 100 and
makes ten investments of 8.4 at times drawn uniformly over the first five years. Each investment is
held for a period drawn uniformly between two and seven years. With probability 12\% it is written
off; otherwise its multiple of cost is $\exp(0.50 + 0.75 Z)$ with $Z$ standard normal, so the mean
multiple over all deals is $0.88\,e^{0.50 + 0.75^2/2} = \mcMeanMult$. Fees are 2\% of commitments for
five years and 1.5\% of the cost of unrealised deals for the next five. Deals are independent within and across funds, and every investment is sold
at the end of its holding period. Ten investments of 8.4 use 84 of the commitment. Fees are 10 in the first five years and at most
$5\times0.015\times84 = 6.3$ in the next five, so paid-in capital could reach 100.3 if every deal were
still unrealised in each of years 5 to 9. The simulation therefore enforces the commitment: fee calls stop once paid-in capital would pass 100.
The limit binds in \mcCapBindN{} of the 6,000 funds, and the largest paid-in capital is \mcPaidMax. The draws come from
NumPy's \texttt{default\_rng} with seed 203. The parameters are illustrative, chosen to give a wide
spread of outcomes; they are not estimates for any market.

\Cref{tab2:mc} summarises the results over all \mcN{} funds. The medians of the three IRRs are taken
separately; the fee and carry wedges are medians of the fund by fund differences, so they need not
equal the differences of the medians. The median fee and carry wedges add up to about five percentage
points of IRR, and the waterfall rule decides how often the GP is paid at all.

\begin{table}[t]
  \centering
  \caption{Simulated funds: \mcN{} funds with commitments of 100. Amounts are per 100
  committed. Carry received is before any clawback; means are over all funds; IRR costs are relative
  to the whole-fund waterfall.}
  \label{tab2:mc}
  \small
  \begin{tabular}{@{}lr@{}}
    \toprule
    Median deal level IRR before fees (\%) & \mcPreFeeIrr\\
    Median IRR after fees (\%) & \mcGrossIrr\\
    Median IRR after fees and carry (\%) & \mcNetIrr\\
    Median fund by fund wedge: fees, carry (percentage points) & \mcWedgeFees, \mcWedgeCarry\\
    Median multiple of paid-in capital before and after carry & \mcTVPI, \mcNetTVPI\\
    \bottomrule
  \end{tabular}

  \medskip
  \begin{tabular}{@{}lrrr@{}}
    \toprule
     & Whole & Deal & Deal by deal,\\
     & fund & by deal & netted\\
    \midrule
    Funds paying no carry (\%) & \mcZeroWf & \mcZeroDbd & \mcZeroPool\\
    Mean carry received & \mcWfPaid & \mcDbdCarry & \mcPoolCarry\\
    Funds owing a clawback (\%) & \mcWfOverProb & \mcClawProb & \mcPoolClawProb\\
    Mean clawback owed & \mcWfOverAll & \mcOwed & \mcPoolOwed\\
    95th percentile of the clawback & \mcWfOverNinetyFive & \mcOwedNinetyFive & \mcPoolOwedNinetyFive\\
    LPs' median IRR cost after clawback (pp) & base & \mcDbdIrrCost & \mcPoolIrrCost\\
    \bottomrule
  \end{tabular}
\end{table}

Under the whole-fund waterfall \mcZeroWf\% of funds pay no carry, because they never clear the hurdle
after fees; under the pure deal-by-deal waterfall almost every fund pays some, because one
profitable deal is enough, and that is what drives the clawback rates in \cref{tab2:mc}.

Even the whole-fund waterfall can pay too much. When capital is called for fees after the catch-up is
complete, the profit on which the carry was computed later shrinks, and the split tier never gives
the excess back (\cref{prop2:catchup}). In \mcWfOverProb\% of the simulated
funds the whole-fund carry paid exceeds 20\% of final profit, by \mcWfOverMean{} on average when it
does, which is why a whole-fund waterfall can still need a final reconciliation of carry. The entitlement used in \cref{tab2:mc} is therefore $\min\{G^{\mathrm{seq}}, c\,(D_{\mathrm{tot}}-P_{\mathrm{tot}})^{+}\}$,
where $G^{\mathrm{seq}}$ is the carry paid by the sequential waterfall, $P_{\mathrm{tot}}$ is total paid-in
capital and $D_{\mathrm{tot}}-P_{\mathrm{tot}}$ is the final profit. The minimum is the model's final
recalculation of carry, a terminal profit true-up in the simulation that need not equal what a fund's
actual clawback clause would recover. Every clawback in \cref{tab2:mc} is measured against this
entitlement, not against the full recomputation of \cref{def2:clawback}. Carry received less the mean clawback is therefore
\mcWfCarry{} to two decimals in every column of \cref{tab2:mc}, because each rule ends at this
entitlement and assumes the GP repays in full, apart from the rare funds where a deal-level rule pays
less than it and nothing is owed. A contractual clawback may differ, since it can be capped net of tax,
and recovering it depends on the GP being able to pay.

The clawback rates in \cref{tab2:mc} are high because of the model. The pure deal level rule measures profit before fees (fees average \mcFeesMean{} per 100 committed) and ignores losing deals, and the netted rule counts realised losses but still not fees, while about one
deal in eight is written off and the deals are independent, so nearly every fund has a shortfall
between the two rules. Under these distribution and true-up rules a simulated fund can owe back carry it has already
received; the figures do not estimate how often real funds owe clawbacks.

The more instructive result is the spike in panel (a) of \cref{fig2:mc} for the netted case: in
\mcPoolFeeCluster\% of funds the clawback is exactly 20\% of the fees paid. Those are successful funds
where both rules give the GP 20\% of profit, but the deal level rule measures profit before fees and
the whole-fund rule after them. Even when the clawback is paid in full the LPs lose, because carry is
paid early and returned late without interest. Panel (b) shows the gross to net mapping fund by fund.

\begin{figure}[t]
  \centering
  \includegraphics[width=\textwidth]{ch02/fig06_simulated_funds}
  \caption{Simulated funds, amounts per 100 committed. (a) Distribution of the clawback owed at the end of
  the fund under the pure deal-by-deal rule and with netting of realised losses. (b) Deal level IRR before fees against the LPs' net IRR under
  the whole-fund waterfall, for a random half of the simulated funds so that the points stay legible.}
  \label{fig2:mc}
\end{figure}

\FloatBarrier
% =====================================================================================
\section{How do hedge fund fees accrue?}
\label{sec2:hf}

\paragraph{Conventions}
A hedge fund's fees are computed each year (or each quarter) on its NAV. Let $A_{t-1}$ be the NAV at the
start of year $t$, $g_t$ the gross return, $m$ the management fee rate, $p$ the incentive fee rate and
$M_{t-1}$ the high-water mark, the highest NAV after fees at any earlier year end (initially $A_0$). This
chapter uses the following order, which is common but not universal: the management fee is $m A_{t-1}$,
deducted at the year end; the incentive fee is charged on the NAV after the management fee,
$P_t = A_{t-1}(1+g_t) - mA_{t-1}$; and with a hurdle rate $h$ the threshold for the incentive fee is
$M_{t-1}(1+h)$. The high-water mark is set after all fees, and the examples have no subscriptions or
redemptions during the year. Then
\begin{equation}\label{eq2:hffee}
  I_t = \begin{cases}
    p\,\bigl(P_t - M_{t-1}(1+h)\bigr)^{+} & \text{hard hurdle},\\[2pt]
    p\,(P_t - M_{t-1})\,\mathbf{1}\{P_t > M_{t-1}(1+h)\} & \text{soft hurdle},
  \end{cases}
\end{equation}
and then $A_t = P_t - I_t$ and $M_t = \max\{M_{t-1}, A_t\}$.
With $h = 0$ the two cases agree. In both chapters on fees a soft hurdle means that once the fee is
earned it is a share of the whole gain, and a hard hurdle that it is a share of the gain above the
threshold only. The soft hurdle here pays on the whole gain as soon as the threshold is passed, a
catch-up with no range; Chapter~9 writes it with a full catch-up, under which the manager takes all of
the gain above the threshold until the fee is $p$ of the whole gain. The two give the same fee once the
return is above $h/(1-p)$. Without a high-water mark, $M_{t-1}$ is replaced by $A_{t-1}$, so the
manager is paid on every year's gain whatever happened before.

\begin{example}{Four years of hedge fund fees}\label{ex2:hf}
A fund starts at 100 and charges 2\% and 20\% with a high-water mark and no hurdle. Gross returns are
15\%, $-10\%$, 8\% and 20\%. \Cref{tab2:hf} shows the calculation. In year 1 the NAV after the
management fee is \hfPreOne, so the incentive fee is $0.2\times(\hfPreOne-100) = \hfIncOne$ and the mark
rises to \hfHwmOne. Years 2 and 3 end below the mark and pay no incentive fee, although year 3 is
profitable. In year 4 the NAV before the incentive fee reaches \hfPreFour, and the fee is charged only
on the gain above \hfHwmThree. Total fees are \hfTotalFees. The gross value would be \hfGrossEnd, a
compound return of \hfGrossCagr\% a year; the investor earns \hfNetCagr\%. Without the high-water mark
the manager would also have earned \hfNoHwmIncThree{} in year 3, and total fees would be
\hfNoHwmTotal. A 5\% hard hurdle cuts the year 1 fee to $0.2\times(\hfPreOne-105) = \hfHardIncOne$; a
5\% soft hurdle leaves it at \hfSoftIncOne{} because the threshold is passed.
\end{example}

\begin{table}[t]
  \centering
  \caption{Fees of \cref{ex2:hf}: 2\% management fee and 20\% incentive fee with a high-water mark.}
  \label{tab2:hf}
  \begin{tabular}{lrrrr}
    \toprule
     & Year 1 & Year 2 & Year 3 & Year 4 \\
    \midrule
    Opening NAV $A_{t-1}$         & \hfOpenOne & \hfOpenTwo & \hfOpenThree & \hfOpenFour \\
    Gross NAV                      & \hfGrossOne & \hfGrossTwo & \hfGrossThree & \hfGrossFour \\
    Management fee                 & \hfMgmtOne & \hfMgmtTwo & \hfMgmtThree & \hfMgmtFour \\
    NAV before incentive fee $P_t$ & \hfPreOne & \hfPreTwo & \hfPreThree & \hfPreFour \\
    Incentive fee $I_t$            & \hfIncOne & \hfIncTwo & \hfIncThree & \hfIncFour \\
    Closing NAV $A_t$              & \hfCloseOne & \hfCloseTwo & \hfCloseThree & \hfCloseFour \\
    High-water mark $M_t$          & \hfHwmOne & \hfHwmTwo & \hfHwmThree & \hfHwmFour \\
    Net return (\%)                & \hfNetOne & \hfNetTwo & \hfNetThree & \hfNetFour \\
    \bottomrule
  \end{tabular}
\end{table}

\paragraph{Modified high-water marks}
After a large loss a traditional high-water mark can leave the manager with no realistic prospect of an
incentive fee for years, which gives it a reason to take more risk or to close the fund and start again.
A \emph{modified} high-water mark, negotiated occasionally after a loss, has two parameters. After a loss of
$\Lambda = M - A_{\min}$ from the mark $M$ to the trough $A_{\min}$, the manager earns a reduced rate
$p' < p$ on gains above the latest NAV after fees until the NAV has recovered a multiple $\mu > 1$ of the
loss, that is until it reaches $A_{\min} + \mu\Lambda$, and the full rate above that level. The investors
bear the cost: during the recovery they pay the reduced rate on gains that a traditional mark would
leave free of incentive fees. Because $\mu > 1$ the target lies above the old mark, so the full rate
returns only once the investors are more than made whole. If the NAV falls below the trough before reaching
the target, the recovery restarts from the new trough, still measured against the old mark $M$, so
$\Lambda$ and the target are recomputed with the new $A_{\min}$.

\begin{example}{A loss and a recovery}\label{ex2:modhwm}
Take returns after management fees of $-25\%$ followed by five years of $+12\%$, a 20\% incentive fee,
and a modified mark with $p' = 10\%$ and $\mu = 1.25$. These values are illustrative; halving the rate and
asking for a quarter more than the loss keeps the arithmetic easy to follow, and any $p' < p$ and
$\mu > 1$ give the same qualitative picture. The loss is 25 from a mark of 100, so the target is $75 + 1.25\times 25 = \modTarget$. Under the traditional mark the first incentive fee
is \modTradIncFour{} in year 4; under the modified mark the manager earns \modModIncTwo{} in year 2,
\modModIncThree{} in year 3 and \modModIncFour{} in year 4. In year 5 the NAV passes the target, so the fee
is $0.1$ times the gain up to \modTarget{} plus $0.2$ times the gain above it, \modModIncFive. Over the
six years the manager earns \modModSum{} under the modified mark and \modTradSum{} under the traditional
one, and the investor's NAV ends at \modModCloseSix{} instead of \modTradCloseSix. The investor pays more in fees during the recovery (\cref{fig2:hwm}). Whether
that buys a manager less inclined to gamble after the loss depends on the manager's other
incentives, which the model does not capture.
\end{example}

\begin{figure}[t]
  \centering
  \includegraphics[width=\textwidth]{ch02/fig07_hwm}
  \caption{\Cref{ex2:modhwm}. (a) NAV after fees under a traditional and a modified high-water mark
  after a 25\% loss. (b) Incentive fees by year.}
  \label{fig2:hwm}
\end{figure}

\FloatBarrier
% =====================================================================================
\section{What is a fee stream worth to the manager?}
\label{sec2:views}

\subsection{If the fund lasts}

A management company is worth, in large part, the fees it expects to collect while its investors stay. \Cref{ex2:annuity} below values ten years of fees on a fund of 1,000 at \annPV, almost eighteen times the first year's
incentive fee of \annIncOne, so protecting the stream can matter more than any one year's incentive fee.

\begin{proposition}[Fees as a growing annuity]\label{prop2:annuity}
Suppose a fund earns a constant, known gross return $g > m$ each year, above its high-water mark,
survives for $T$ years, and pays fees at each year end, and that some fee is charged, so that
$\phi = m + p(g-m) > 0$. Then the fee in year $t$ is $\phi A_{t-1}$, the NAV grows at $G = (1-p)(g-m)$, and, for a discount rate $d$ with $1+G>0$ and
$1+d>0$, the present value of $T$ years of fees is
\begin{equation}\label{eq2:annuity}
  \mathrm{PV}_T = \frac{\phi A_0}{1+d}\,\frac{1-q^{T}}{1-q}, \qquad q = \frac{1+G}{1+d}\ne 1 .
\end{equation}
If $G<d$ the perpetuity $\lim_{T\to\infty}\mathrm{PV}_T$ exists, and if the fees are discounted at $d = g$
it equals the initial NAV, $A_0$.
\end{proposition}
\begin{proof}
By \cref{eq2:hffee} with $h=0$ and $P_t > M_{t-1} = A_{t-1}$, the fees are
$mA_{t-1} + p\bigl(A_{t-1}(1+g) - mA_{t-1} - A_{t-1}\bigr) = A_{t-1}[m + p(g-m)] = \phi A_{t-1}$, and
$A_t = A_{t-1}(1+g) - \phi A_{t-1} = A_{t-1}(1+G)$ because $g - \phi = (1-p)(g-m) = G$. So
$A_{t-1} = A_0(1+G)^{t-1}$ and
\begin{equation*}
  \mathrm{PV}_T = \sum_{t=1}^{T} \frac{\phi A_0 (1+G)^{t-1}}{(1+d)^{t}}
           = \frac{\phi A_0}{1+d} \sum_{t=0}^{T-1} q^{t}
           = \frac{\phi A_0}{1+d}\,\frac{1-q^T}{1-q}.
\end{equation*}
If $g = d$ and $\phi > 0$, then $G = d - \phi < d$, so $0 < q < 1$ and $1 - q = \phi/(1+d)$. Hence
$\mathrm{PV}_T = A_0(1-q^T)$ for every $T$, and the limit is $A_0$.
\end{proof}

If the fund earns exactly the rate at which its fees are discounted, the fees and the investors'
claim, $A_0(1+G)^T/(1+d)^T$, add up to $A_0$ in present value for every $T$: they are two shares of
one asset, and as $T$ grows the fees take all of it. The result holds only in this deterministic model. The manager's
discount rate need not be the investors' required return, and finite lives and redemptions stop it
happening in practice. What it shows is how much a long-lived franchise can be worth to its manager.

\begin{example}{The value of a fee stream}\label{ex2:annuity}
A fund of 1,000 earns 10\% gross, charges 2\% and 20\%, and the manager discounts at 10\%. Then
$\phi = \annPhi\%$ and the NAV grows at \annNetGrowth\% a year. Ten years of fees are worth \annPV, of which
\annPVm{} is management fees (\annMgmtShare\%) and \annPVi{} incentive fees; twenty years are worth
\annPVtwenty, and the perpetuity is worth \annPVperp, the opening NAV, under the perpetual assumptions. Panel (b) of \cref{fig2:annuity}
shows that the management fee dominates the value unless gross returns are high.
\end{example}

\begin{figure}[t]
  \centering
  \includegraphics[width=\textwidth]{ch02/fig09_annuity}
  \caption{The annuity view of fees for a fund of 1,000 discounted at 10\%. (a) Present value of fees
  against the number of years the fund survives, at a 10\% gross return; the dotted line is a fund
  charging the 2\% management fee with no incentive fee. (b) Present value of ten years of fees by
  component against the gross return.}
  \label{fig2:annuity}
\end{figure}

\subsection{One year at a time}

Over one year the incentive fee is $p\,(P_t - M_{t-1})^{+}$, which has the payoff of a call option on
the fund's NAV struck at the high-water mark. The terminal payoffs coincide, but the fee cannot be traded,
hedged or sold separately from the fund, and the manager can influence the NAV it is struck on. The
Black-Scholes price is still a useful benchmark, as in \cref{sec2:options}. Because the management fee is charged on the opening NAV $A_{t-1}$, the fee is exactly $p$ calls on the
gross NAV struck at $M_{t-1} + mA_{t-1}$. Treating the fee instead as a call on the NAV net of the
management fee, $S = (1-m)A_{t-1}$, struck at the mark is a pricing approximation and not an identity:
the fee is charged on opening assets, so it is not removed in proportion from the terminal payoff, and
the two strikes give different values. With the inputs used below and a 2\% management fee, the
exact form is worth \optExactAtSigTwenty{} per 100 of mark when $S$ is at the mark, against
\optAtSigTwenty{} for the approximation, so little is lost, and the approximation puts NAV and mark on
one axis. The approximation is
\begin{equation}\label{eq2:hfoption}
  \mathcal{I}_0 \approx p\, C(S, M, 1, r, \sigma).
\end{equation}
\textcite{goetzmann2003} value the full multi-year high-water mark contract in a continuous time
model and find that the fees are a claim on a significant share of the investors' wealth.

\begin{proposition}[Where the incentive to add risk is strongest]\label{prop2:vegapeak}
As a function of $S$, the vega of \cref{eq2:hfoption}, $p\,S\,\varphi(d_1)\sqrt{T}$, is largest at
\begin{equation}\label{eq2:vegapeak}
  S^{*} = M\exp\bigl\{(\tfrac12\sigma^2 - r)\,T\bigr\},
\end{equation}
close to the high-water mark when $\sigma^2/2$ and $r$ are small.
\end{proposition}
\begin{proof}
Up to the constant factor $p\sqrt T$ the vega is
\begin{equation*}
  \nu(S) = S\varphi(d_1), \qquad d_1 = \frac{\ln(S/M) + (r+\sigma^2/2)T}{\sigma\sqrt T},
  \qquad \frac{\partial d_1}{\partial S} = \frac{1}{S\sigma\sqrt T}.
\end{equation*} Since $\varphi'(x) = -x\varphi(x)$,
\begin{equation*}
  \nu'(S) = \varphi(d_1) + S\,\varphi'(d_1)\,\frac{1}{S\sigma\sqrt T}
          = \varphi(d_1)\Bigl(1 - \frac{d_1}{\sigma\sqrt T}\Bigr),
\end{equation*}
which is positive for $d_1 < \sigma\sqrt T$ and negative above, so the maximum is at $d_1 = \sigma\sqrt T$,
that is $\ln(S/M) + (r + \sigma^2/2)T = \sigma^2 T$, which gives \cref{eq2:vegapeak}.
\end{proof}

With a 20\% incentive fee, $T=1$, a 3\% rate and 20\% volatility, the one-year fee is worth \optAtSigTwenty{} per 100 of high-water mark when $S$, the NAV net of the management fee, is at the mark,
\optAtSigTen{} at 10\% volatility, \optUpSigTwenty{} when the NAV is 30\% above the mark, and only
\optDownSigTwenty{} when it is 20\% below. A manager 20\% below the mark who doubles volatility to 40\%
raises the value to \optDownSigForty. Per point of volatility the fee gains \optVegaAt{} at the mark,
\optVegaDown{} at 20\% below and \optVegaUp{} at 30\% above (panel (b) of \cref{fig2:optionview}). The
pressure to add risk is strongest near the mark and fades far above it, the same pattern as
\cref{eq2:vega} for private equity carry.

The option view also explains why, in the models here, fees take a larger share of expected gains from
riskier funds. The comparison holds the expected gross return fixed and changes only the dispersion
around it, so it is a property of the fee formula and makes no claim about how real funds' returns
and risks are related.

\begin{proposition}[Convexity of the incentive fee]\label{prop2:convex}
If the gross gain above the mark is $X \sim \N(\mu, s^2)$, the expected incentive fee
$p\,\E[X^{+}] = p\,[\mu\Phi(\mu/s) + s\varphi(\mu/s)]$ is strictly increasing in $s$, with
$\partial\,\E[X^{+}]/\partial s = \varphi(\mu/s)$.
\end{proposition}
\begin{proof}
Write $X = \mu + sZ$. Then $\E[X^{+}] = \int_{-\mu/s}^{\infty}(\mu + s z)\varphi(z)\dd z = \mu\Phi(\mu/s) + s\varphi(\mu/s)$,
using $\int_a^\infty z\varphi(z)\dd z = \varphi(a)$ and the symmetry of $\varphi$. Differentiate with $u = \mu/s$,
$\partial u/\partial s = -\mu/s^2$ and $\varphi'(u) = -u\varphi(u)$:
\begin{align*}
  \frac{\partial}{\partial s}\E[X^{+}]
  &= \mu\varphi(u)\Bigl(-\frac{\mu}{s^2}\Bigr) + \varphi(u) + s\bigl(-u\varphi(u)\bigr)\Bigl(-\frac{\mu}{s^2}\Bigr)\\
  &= -\frac{\mu^2}{s^2}\varphi(u) + \varphi(u) + \frac{\mu^2}{s^2}\varphi(u) = \varphi(u) > 0. \qedhere
\end{align*}
\end{proof}

Panel (c) of \cref{fig2:optionview} shows the effect over ten years, using \shareN{} simulated paths
(seed 204). Each year's gross return is lognormal, $1+g = \exp\{\ln 1.08 - s^2/2 + sZ\}$, so the
expected gross return is 8\% a year whatever the volatility $s$; fees are 2\% and 20\%, crystallised at
each year end, with the mark updated after fees. Fees are measured as expected gross minus expected
net terminal value, as a share of the expected gross gain. With no volatility the fees take \shareDet\% of
the expected gross gain. With a high-water mark they take \shareHwmTen\% at 10\% volatility,
\shareHwmTwenty\% at 20\% and \shareHwmForty\% at 40\%. Without a high-water mark the shares are
\shareNoHwmTen\%, \shareNoHwmTwenty\% and \shareNoHwmForty\%. The high-water mark is worth most to
investors in the funds where the incentive fee is most valuable to the manager.

\begin{figure}[t]
  \centering
  \includegraphics[width=\textwidth]{ch02/fig08_option_view}
  \caption{The option view of hedge fund fees (20\% incentive fee, 3\% rate). (a) Black-Scholes model value
  of one year's incentive fee per 100 of high-water mark against the ratio of $S$, the NAV net of the management fee, to the mark. (b) Its sensitivity to
  one point of volatility. (c) Simulated fees as a share of the expected ten year gross gain, over \shareN{} paths; the
  dotted line is the case of no volatility.}
  \label{fig2:optionview}
\end{figure}

\FloatBarrier
% =====================================================================================
\section{What do the fee terms encourage?}
\label{sec2:behaviour}

A payoff shows what a manager stands to gain. That suggests what to watch for in its conduct, but it
does not show what the manager will do. In the normal model the incentive fee always gains from
volatility (\cref{prop2:convex}), most strongly near the mark (\cref{prop2:vegapeak}), while carry,
by \cref{eq2:vega}, can lose value far above the catch-up threshold.

A GP or hedge fund manager well below its threshold holds an out-of-the-money claim, which gains
in value from more volatility while the investors bear the extra risk. That makes risk more
attractive to the manager without settling what it will choose. Its own capital at risk, the value of
future fees and the threat of withdrawals can outweigh that pull.

A manager with large gains early in a fee period is in a different position. Much of the fee is
already accrued, and a reversal before it crystallises could reduce or remove it, so a risk-averse
manager may cut risk for the rest of the period to lock it in, even when the investors would prefer
it to keep investing, though the fee's option value still rises with volatility.

When the management fee is most of the value of the fee stream (\cref{ex2:annuity}), raising assets
can pay the manager more than earning returns, which creates an incentive to grow the fund, possibly
beyond the capacity of its strategy. A manager may then keep growing it until the lost returns, the weaker record and the redemptions
that follow cost it more than the extra fees bring in. How strong this is depends on the fee base: a fee on NAV rewards growth in the fund's value,
while a fee on gross assets also rewards leverage. A manager can also protect the fee stream by
holding a portfolio close to its benchmark while charging active fees, so that the investors pay for
active management and receive an index.

Where positions are valued with discretion, a manager near a threshold has a reason to report values
that clear it or to smooth reported returns, which Chapter~1 showed lowers measured volatility; having
a reason does not mean it acts on it. Exit
timing in private funds is a separate matter. Whether a GP near the hurdle gains from delaying a sale
or from bringing one forward depends on whether the fund's value already clears the hurdle balance and on
how fast that balance compounds against the asset's expected growth; a GP far below the hurdle may have a
weaker immediate carry incentive, although future fees, reputation and its own invested capital can still
reward effort.

\paragraph{Evidence}
The evidence supports some of these predictions and qualifies others, and it is correlational. \textcite{agarwal2009} find that
hedge funds with higher incentive fees, high-water marks and larger manager ownership deliver better
performance, consistent with incentives that align interests. \textcite{brown2001} find that funds
that do well in the first half of a year reduce volatility in the second half and poor performers
raise it, but the shifts track performance relative to other funds more closely than the distance to
the high-water mark, and the risk that a poorly performing fund is closed moderates the incentive to
gamble. Career concerns and the value of the future fee stream (\cref{prop2:annuity}) therefore pull
against the option incentive.

\paragraph{Fees on fees}
A fund of funds\index{Fund of funds} adds its own fees, 1\% and 10\% in the illustration below (terms vary), on top of the
fees of the funds it holds, which are charged first, inside the underlying funds' NAVs.
The two layers are not simply additive, because each underlying manager is paid on its own gains.

\begin{proposition}[Netting risk]\label{prop2:netting}
Let a portfolio hold weights $w_i\ge 0$, $\sum w_i = 1$, in funds whose gains above their thresholds are
$X_i$. The portfolio pays incentive fees $p\sum_i w_i X_i^{+}$, while one fund with the same holdings and
the same fee rate would pay $p\,(\sum_i w_i X_i)^{+}$. Pathwise and in expectation,
\begin{equation}\label{eq2:netting}
  p\sum_i w_i X_i^{+} \;\ge\; p\Bigl(\sum_i w_i X_i\Bigr)^{+},
\end{equation}
and the gap does not shrink under a mean-preserving spread of the $X_i$ under the weights $w_i$,
defined in the proof.
\end{proposition}
\begin{proof}
$x\mapsto x^{+}$ is convex, so by Jensen's inequality for the discrete distribution with weights $w_i$,
$(\sum w_i X_i)^{+} \le \sum w_i X_i^{+}$ on every path; take expectations.
Call $(X_i')$ a \emph{mean-preserving spread}\index{Mean-preserving spread} of $(X_i)$ under the weights $w_i$ if $X_i' = X_i + \varepsilon_i$
where, for each value $x$ taken by the $X_i$, the weighted mean of the $\varepsilon_i$ over the funds with
$X_i = x$ is zero: $\sum_{i:X_i=x} w_i\varepsilon_i = 0$. Equal weighted means alone are not enough; with
two equally weighted funds, $X=(-1,1)$ and $X'=(0,0)$ have the same mean but the gap falls from $\tfrac12$
to zero. Fix a value $x$ and let $W_x = \sum_{i:X_i=x} w_i > 0$. Jensen's inequality for the weights
$w_i/W_x$ over that group gives
\[
  \sum_{i:X_i=x} \frac{w_i}{W_x}(x+\varepsilon_i)^{+} \;\ge\; \Bigl(x + \sum_{i:X_i=x}\frac{w_i}{W_x}\varepsilon_i\Bigr)^{+} = x^{+}.
\]
Multiplying by $W_x$ and summing over the values $x$ gives $\sum w_i X_i'^{+} \ge \sum w_i X_i^{+}$, so
the left side of \cref{eq2:netting} does not fall. Summing $\sum_{i:X_i=x} w_i\varepsilon_i = 0$ over $x$
gives $\sum w_i X_i' = \sum w_i X_i$, so the right side is unchanged and the gap does not shrink.
\end{proof}

With equal amounts in two funds that gain 30\% and lose 10\% before fees, each charging 1.5\% of assets and 20\% of the gain above the management fee\ix{Management fee},
the fees are $1.5 + 0.2\times(30-1.5) = \netFeeA\%$ and \netFeeB\% (no incentive\ix{Incentive compatibility} fee on a loss), an average of \netPaid\% on a gross portfolio return of \netGross\%. A single fund
with the same terms on the combined return would charge \netSingle\%. The investor pays an effective
incentive rate of \netEffRate\% on the portfolio's gain above the management fee. \Cref{fig2:fof}
repeats the comparison for ten funds with normally distributed one-year gross returns\ix{Gross return}, mean 10\% and
volatility 15\%, equal pairwise correlation, and fees charged once at the year end with no high-water
mark (200{,}000 draws, seed 205). With
uncorrelated funds the underlying fees take \fofUnderZero\% of the expected gross return, against
\fofSingleZero\% for a single fund with the same terms, and adding the fund of funds\ix{Fund of funds} layer brings the
total to \fofTotalZero\%. As correlation rises to one the netting gap closes (\fofUnderOne\% against
\fofSingleOne\%), but from one side only: each fund's expected fee does not depend on correlation, so the
underlying fees barely move, and it is the single fund's charge that rises to meet them once all the
funds win or lose together. \textcite{brown2003} study these fees
on fees and their effect on the net returns of funds of funds.

\begin{figure}[t]
  \centering
  \includegraphics[width=0.75\textwidth]{ch02/fig10_fof}
  \caption{Expected fees as a share of expected gross return for ten equally weighted funds (10\%
  expected return, 15\% volatility, 1.5\% and 20\% fees) against their correlation: underlying fees only,
  underlying fees plus a 1\% and 10\% fund of funds layer, and a single fund charging 1.5\% and 20\% on the
  combined return. One-year normal returns with equal pairwise correlation; fees charged once at the
  year end with no high-water mark.}
  \label{fig2:fof}
\end{figure}

\FloatBarrier
% =====================================================================================
\section{What does the contract do about these incentives?}
\label{sec2:terms}

Four conflicts come out of the payoff calculations. Carry is an option, so a GP\ix{General partner} below the
hurdle\ix{Hurdle rate} can gain from risk the investors\ix{Investor} bear (\cref{sec2:options,sec2:behaviour}). Carry measured deal
by deal can pay the GP before the fund as a whole has earned anything (\cref{sec2:dbd,sec2:mc}). A GP
far below the hurdle has more to gain from its next fund than from this one (\cref{sec2:behaviour}).
And a fee on assets rewards size whatever the return (\cref{sec2:mgmtfees}). Because taking part in management may put an LP\ix{Limited partner}'s limited liability at risk under the governing law,
an LP usually answers these pressures with rights negotiated in advance and in writing, and leaves
day-to-day control to the GP. \Cref{tab2:terms} lists which term answers which.

The investor learns the strategy, team and terms from a \emph{private placement memorandum}\index{Private placement memorandum} (PPM),
which describes the offer and its risks; the investor's enforceable rights and remedies come mainly
from the governing documents, any representations made and the applicable securities law. It joins by signing a
\emph{subscription agreement}\index{Subscription agreement}, which fixes its commitment\ix{Commitment} and confirms that it is eligible to invest.
Its rights come from the \emph{limited partnership agreement\ix{Limited partnership agreement}}\index{Limited partnership agreement} (LPA), together with any \emph{side letter}\index{Side letter}
it negotiates for itself and, through a \emph{most favoured nation}\index{Most favoured nation} clause, terms granted to comparable
investors, usually those of similar or smaller size. The LPA also fixes what the fund may buy, how
large it may grow and how long it lasts, and the penalties for an investor that defaults\ix{Default} on a call.
Investors do not usually see the operating agreement of the management company, which divides fees
and carry among the principals.

\begin{table}[t]
  \centering
  \caption{Contract terms and the pressure each one answers.}
  \label{tab2:terms}
  \small
  \begin{tabular}{@{}>{\raggedright\arraybackslash}p{0.37\textwidth}>{\raggedright\arraybackslash}p{0.37\textwidth}>{\raggedright\arraybackslash}p{0.18\textwidth}@{}}
    \toprule
    Term & Pressure it answers & Derived in\\
    \midrule
    GP commitment paid in cash & Gains from risk below the hurdle & \cref{sec2:options}\\
    Whole-fund waterfall, clawback, escrow & Carry paid before the fund has earned it & \cref{sec2:dbd,sec2:mc}\\
    Key person clause, successor-fund limits & Attention moving to the next fund & \cref{sec2:behaviour}\\
    Fee on invested capital after the investment period & A fee that rewards size & \cref{sec2:mgmtfees}\\
    Allocation rules, fee offsets, advisory committee & Conflicts between funds and with portfolio companies & this section\\
    Removal for cause or without cause & Failure of the other terms & this section\\
    \bottomrule
  \end{tabular}
\end{table}

\paragraph{Sharing the downside}
The direct answer to the option incentive\ix{Incentive compatibility} is to make the GP lose money when the investors do. A
commitment by the GP gives the principals a share of investment losses alongside the LPs, and
practitioners call it \emph{hurt money}\index{Hurt money} because it should be large next to the principals' own wealth.
Its size relative to that wealth, how it is financed and how it is shared among the principals say
more about alignment than the headline percentage. A commitment funded by waiving management fees\ix{Management fee}
puts at risk fee income that the LPs paid for and leaves the principals' after-tax savings untouched,
which is why many investors give it less weight than cash.
The whole-fund\ix{Distribution waterfall!European} waterfall and the clawback\ix{Clawback} deal with timing: carry waits until the LPs have their
capital and preferred return, and carry paid too early comes back at the end. \Cref{sec2:mc} showed
that late repayment still costs the LPs the time value of the carry, and a clawback is worth only what
the GP can pay, which is why agreements add escrow, interim clawback tests or guarantees from the
principals.

\paragraph{Keeping attention on the fund}
Investors back people, and a manager far below its hurdle may have weaker reasons to spend its time
on the existing portfolio, especially when it can raise another fund. A
\emph{key person\ix{Key person clause}}\index{Key person} clause names the principals the fund depends on and requires each to give it a
stated share of working time. A typical trigger is that two of three named principals leave or stop
meeting that test. Many agreements then suspend the investment period, which ends unless investors
vote to reinstate it, usually after the GP proposes replacements; others end it at once. Limits on
raising a successor fund\ix{Successor fund} until most of the current one is invested answer the same drift directly.

\paragraph{Conflicts between funds}
A manager with several funds must decide which of them gets a deal, may offer co-investment to some
LPs and not others, may sell an asset from one fund to another and may collect transaction and
monitoring fees\ix{Transaction fee} from portfolio companies. The LPA handles these with rules set in advance: an
allocation policy, an offset of portfolio company fees against the management fee, limits on
recycling\ix{Recycling} proceeds into new deals and limits on fund level borrowing and guarantees. Cases the rules
do not settle go to the \emph{limited partner advisory committee}\index{Limited partner advisory committee} (LPAC), made up of representatives
of the larger investors. It reviews conflicts and, where the agreement gives it that authority,
valuation questions and waivers of the investment restrictions. It does not take investment
decisions, since that would compromise the LPs' passive\ix{Index fund} status. The Institutional Limited Partners
Association, a body of private fund investors, publishes principles on alignment, governance and
transparency that many investors use as a reference when they negotiate \parencite{ilpa2019}.

\paragraph{Replacing the manager}
When the rules fail, the investors' last protection is to replace the GP, and agreements usually
offer two routes with thresholds and consequences negotiated fund by fund. Removal \emph{for cause}
follows a defined wrong such as fraud or a material breach of the agreement, often as found by a court
or an arbitrator. Since the GP is at fault the vote needed tends to be lower, and the GP can lose some
or all of its carry on unrealised investments and its future fees. Removal \emph{without cause} needs
only lost confidence, and it asks for more, commonly two thirds to three quarters of interests,
because the investors are taking away fees and carry the GP expected to earn through no fault of its
own. The GP then usually keeps the carry earned to date and may receive compensation such as a period
of fees. Removing the GP, ending the investment period and dissolving the fund are different outcomes:
the first changes who manages the fund, the second stops new investments while the existing ones are
managed to exit, and the third winds up the partnership. An agreement may attach a different vote to
each.

The same logic applies inside the management company to principals who leave. Carry usually vests
over the investment period. A \emph{good leaver}\index{Good leaver}, who leaves for an accepted reason such as illness or
retirement\ix{Retirement}, keeps what has vested; a \emph{bad leaver}\index{Bad leaver}, who joins a competitor or leaves after
misconduct, forfeits some or all of it. Vesting, forfeiture, repurchase and the definition of cause all depend on
the management company's own agreements.

\FloatBarrier
% =====================================================================================
\section{How do property sponsors share profits?}
\label{sec2:promote}

Real estate syndications and joint ventures\ix{Venture capital} share profits between an investor and a sponsor that
finds, finances and manages the property. The sponsor usually invests a small share of the equity and
earns a \emph{promote}, a share of distributions above its pro rata\ix{Pro rata} share, in tiers set by the
investor's return. The tiers are often defined by IRR\ix{Internal rate of return} hurdles\ix{Hurdle rate} and sometimes by equity multiples\ix{Comparable companies}\ix{Equity multiple}\ix{Valuation multiple}.

\begin{definition}[Tiered promote]\label{def2:promote}
Let $C_i$ be the total equity contributed at time $t_i$, of which the investor provides the share $s$,
that is $sC_i$, and the sponsor $1-s$. Tier $j$ has a promote
$\pi_j$ and a hurdle: an IRR $h_j$ or an equity multiple $\mu_j$, measured on the investor's own cash
flows\ix{Cash flow} at their actual dates. Within tier $j$ the sponsor first takes $\pi_j$ of each dollar
distributed and the rest is split pro rata, so the investor receives $(1-\pi_j)s$ per dollar; as a
modelling assumption that follows a common contractual form, the promote applies to distribution
dollars, not to profit dollars. Tier $j$ begins where tier $j-1$ ends
and ends when the investor's flows reach hurdle $j$, measured by the investor balance
\begin{equation}\label{eq2:tierbalance}
\begin{aligned}
  B_j(t) &= \sum_{i:\,t_i\le t} s\,C_i (1+h_j)^{t-t_i} - \sum_{u<t} L_u (1+h_j)^{t-u} - \ell_{<j}(t)
  \quad\text{or}\\
  B_j(t) &= \mu_j \sum_{i:\,t_i\le t} s\,C_i - \sum_{u<t} L_u - \ell_{<j}(t),
\end{aligned}
\end{equation}
where $L_u$ are the investor's distributions at earlier dates and $\ell_{<j}(t)$ is what the investor has
already received from tiers $1,\dots,j-1$ of the distribution at $t$. A distribution $X$ at $t$ fills
tier $j$ with $x_j = \min\{X_j, B_j^{+}/((1-\pi_j)s)\}$, where $X_j$ is what is left of $X$ after the
earlier tiers.
\end{definition}

The size of each tier follows from the per dollar split. If the investor still needs $B_j>0$ after
what earlier tiers paid it at this date, and it receives $(1-\pi_j)s$ per dollar, the tier holds
$B_j/((1-\pi_j)s)$ dollars of the distribution before anything enters the next tier. Without the term
$\ell_{<j}(t)$ the tiers would overlap within one distribution and count the same dollars twice. The
last tier has no hurdle.

\begin{example}{A five year property}\label{ex2:promote}
An investor provides 9 and a sponsor 1 of the equity of 10 for a property. It pays cash flow of 0.55 in
year 1, rising 2\% a year, and is sold in year 5 for net proceeds of \reP, paid together with that year's cash flow. Tier 1 is pro rata until
the investor's IRR reaches 8\%, tier 2 carries a 20\% promote until the investor's IRR reaches 12\%, and
tier 3 a 30\% promote. Total distributions are \reDist, the project IRR is \reProjIrr\% and the profit is
\reProfit. The distributions of years 1 to 4 are all in tier 1. In year 5 the investor's 8\% balance is
$9\times1.08^5$ less its earlier receipts compounded at 8\%, which is \reBalOne, so tier 1 takes
$\reBalOne/0.9 = \reTierOneSize$ of the \reYrFive{} paid that year. The 12\% balance net of that tier 1
receipt is \reBalTwo, so tier 2 takes $\reBalTwo/(0.8\times0.9) = \reTierTwoSize$, and the remaining
\reTierThreeSize{} goes to tier 3. In all the investor receives \reTierOneInv{} in tier 1, \reTierTwoInv{} in tier 2 and \reTierThreeInv{}
in tier 3, a total of \reInvTot{} and an IRR of \reInvIrr\%. The sponsor receives \reTierOneSpo,
\reTierTwoSpo{} and \reTierThreeSpo, a total of \reSpoTot, of which \rePromote{} is promote above its 10\%
pro rata share. The sponsor's IRR is \reSpoIrr\% and it takes \reSpoShareProfit\% of the profit for 10\%
of the equity.
\end{example}

\paragraph{IRR or multiple hurdles}
With equity multiple hurdles of 1.0 and 1.6 times in place of the 8\% and 12\% IRR hurdles, the sponsor
receives \reMulSpoTot{} in \cref{ex2:promote} and the investor's IRR falls to \reMulInvIrr\%. The two
kinds of hurdle respond differently to time. If the cash is spread over ten years instead of five, with the 0.55
cash flow growing at 2\% for ten years and the sale price reduced so that total cash stays \reDist, the project IRR falls to \reTenIrr\%, and so does the investor's, which never reaches the 8\% hurdle; with
IRR hurdles the sponsor gets only its pro rata \reTenSpoIrrTiers{} and no promote, while with multiple
hurdles it still gets \reTenSpoMulTiers, because multiples ignore time. Total cash is the same in both cases, so the comparison isolates timing. An IRR hurdle rewards
speed and can be cleared on a quick sale with a small profit; a multiple hurdle rewards the size of the
profit and can be cleared by holding a mediocre asset for a long time. Many agreements use both, and
require both to be met. \Cref{fig2:promote} shows the sponsor's share of profit against the project IRR
for both designs.

\begin{figure}[t]
  \centering
  \includegraphics[width=\textwidth]{ch02/fig11_promote}
  \caption{The promote of \cref{ex2:promote}. (a) Sponsor's share of total profit against the project
  IRR, varying the sale price, for IRR hurdles and for equity multiple hurdles. (b) Distributions by tier
  in the base case with IRR hurdles.}
  \label{fig2:promote}
\end{figure}

\FloatBarrier
% =====================================================================================
\section{How much should an investor commit each year?}
\label{sec2:pacing}

An investor that wants a stable allocation to private funds must commit before it knows when capital
will be called or returned. \textcite{takahashi2002} give a deterministic model of one fund's cash flows\ix{Cash flow}
that is widely used for this planning. The version here is simplified: the published model also has a
yield term, omitted here, and the cap $\min\{1,\cdot\}$ on the distribution rate is a choice made for this
book. Within each year the opening NAV\ix{Net asset value} grows at $G$, the distribution is taken from the grown NAV,
and the year's contribution is added afterwards; $RC_t$ is a fraction of the uncalled commitment\ix{Commitment}\ix{Uncalled commitment} and
$RD_t$ a fraction of the grown NAV before the contribution.

\begin{definition}[Takahashi-Alexander model]\label{def2:ta}
For a commitment of 1 made at time 0, with contribution rates\ix{Contribution rate} $RC_t$, a growth rate\ix{Growth rate} $G$, a fund life\ix{Fund life} $L$
and a bow $B$, define for $t = 1, 2, \dots$
\begin{align}
  C_t &= RC_t\,U_{t-1}, \qquad U_t = U_{t-1} - C_t, \qquad U_0 = 1, \label{eq2:taC}\\
  D_t &= RD_t\,\mathrm{NAV}_{t-1}(1+G), \qquad RD_t = \min\{1, (t/L)^{B}\}, \label{eq2:taD}\\
  \mathrm{NAV}_t &= \mathrm{NAV}_{t-1}(1+G) + C_t - D_t, \qquad \mathrm{NAV}_0 = 0. \label{eq2:taNAV}
\end{align}
\end{definition}

The uncalled commitment has the closed form $U_t = \prod_{s=1}^{t}(1-RC_s)$, by induction on
\cref{eq2:taC}. A larger bow $B$ keeps distribution rates low for longer, so the NAV builds a higher
hump; $RD_t = 1$ for $t\ge L$ means that from age $L$ the fund pays out everything it holds, so
its NAV after that is only the capital called in that year, $\mathrm{NAV}_t = C_t$. That decays
geometrically when the contribution rate stays constant and below 100\%; in general it follows the
remaining contribution schedule.

\begin{example}{One fund}\label{ex2:ta}
Take $RC = 25\%$, $1/3$ and then 50\% a year, $G = 12\%$, $B = 2.5$ and $L = 12$. By \cref{eq2:taC},
$U_1 = 0.75$, $U_2 = 0.5$ and $U_t = 0.25\times 0.5^{t-3}$ for $t\ge 3$, so \taCalledFive{} of the commitment
is called by year 5 and \taCalledTotal{} in all. The NAV peaks at \taPeakNav{} in year \taPeakYear, total
distributions are \taDistTotal, and the multiple\ix{Valuation multiple} of paid-in capital\ix{Paid-in capital} is \taTVPI. Cumulative net cash\ix{Cumulative net cash flow}
reaches its low of $-\taTroughCash$ in year \taTroughYear{} and turns positive in year \taBreakeven, the
J-curve of Chapter~1. \Cref{fig2:pacingsens} shows that the peak NAV is \taPeakBowLow, \taPeakBowMid{} and
\taPeakBowHigh{} for bows of 1.5, 2.5 and 4, and the multiple is \taTVPIGLow, \taTVPIGMid{} and
\taTVPIGHigh{} for growth rates of 4\%, 12\% and 20\%.
\end{example}

\begin{proposition}[Steady state of a commitment programme]\label{prop2:steady}
Suppose an investor commits 1 every year to funds that follow \cref{def2:ta}, and let $\mathrm{NAV}_a$,
$U_a$, $C_a$ and $D_a$ be the values for one fund at age $a$, with $U_0 = 1$. Assume $G>-1$, every contribution
rate lies in $[0,1]$, and the contribution rate is eventually a constant in $(0,1]$, as in \cref{ex2:ta}. As the programme ages its totals converge to constants:
\begin{equation}\label{eq2:steady}
  \overline{\mathrm{NAV}} = \sum_{a\ge 0}\mathrm{NAV}_a, \qquad \overline{U} = \sum_{a\ge 0} U_a, \qquad
  \overline{D} - \overline{C} = \sum_{a\ge 1} D_a - \sum_{a\ge 1} C_a .
\end{equation}
In particular, steady net cash flow\ix{Net cash flow} per year equals the total net cash flow of one fund over its life.
\end{proposition}
\begin{proof}
The model is linear in the commitment and does not depend on calendar time\ix{Calendar time}, so the fund committed in
year $v$ contributes $\mathrm{NAV}_{t-v}$ to the programme's NAV in year $t$. Summing over $v \le t$ gives
$\sum_{a=0}^{t}\mathrm{NAV}_a$. The same argument applies to $U$, $C$ and $D$. For $a\ge L$, including $a=L$ itself, $RD_a = 1$, so $D_a = \mathrm{NAV}_{a-1}(1+G)$ and
\cref{eq2:taNAV} leaves $\mathrm{NAV}_a = C_a$; and $C_a$ and $U_a$ fall geometrically once the contribution rate is constant (the assumption of the
proposition),
so every sum converges as $t\to\infty$; for $a \ge L$ the year's distribution $D_a$ is the previous NAV
grown at $G$, which decays at the same rate.
\end{proof}

With the parameters of \cref{ex2:ta}, a programme of 1 a year settles at a NAV of \progNavSteady{} and
uncalled commitments of $1 + 0.75 + 0.5 + 0.25\sum_{n\ge 0}0.5^{n} = \progUncalledSteady$, and returns a
net \progNetCash{} a year, equal to the lifetime net cash of one fund. An investor that wants a steady NAV
of 100 must commit \progCommitForHundred{} a year, and its total exposure (NAV plus uncalled
commitments) is then \progOvercommit{} times its NAV (\cref{fig2:pacing}). The ratio sizes the liquidity that must stand behind the uncalled commitments; it is a planning
figure and says nothing by itself about whether a given call can be met.

\begin{figure}[t]
  \centering
  \includegraphics[width=\textwidth]{ch02/fig12_pacing}
  \caption{The Takahashi-Alexander model with $RC = 25\%, 1/3, 50\%, \dots$, $G = 12\%$, $B = 2.5$, $L = 12$.
  (a) Contributions, distributions, NAV and cumulative net cash of one fund. (b) A programme that
  commits 1 every year.}
  \label{fig2:pacing}
\end{figure}

\begin{figure}[!htb]
  \centering
  \includegraphics[width=0.86\textwidth]{ch02/fig13_pacing_sensitivity}
  \caption{Sensitivity of the Takahashi-Alexander model. (a) NAV of one fund for three values of the bow
  $B$. (b) Cumulative net cash for three growth rates $G$.}
  \label{fig2:pacingsens}
\end{figure}

\FloatBarrier
% =====================================================================================
\section{Summary}
\label{sec2:summary}

\begin{itemize}
  \item A limited partnership generally combines limited liability for the LPs with flexible terms and
        pass-through tax, subject to the jurisdiction and the investor; when
        investors need different tax treatment, a master-feeder lets them hold the same portfolio.
  \item An open-end investor can leave only on the fund's terms, which may lock it in, gate it or
        suspend it. A drawdown investor must fund calls it does not control, cannot redeem, and can
        leave only by selling its interest at a negotiated price.
  \item The waterfall pays capital, a compounding preferred return, a catch-up and a split. The
        hurdle balance is nonpositive exactly when the LPs' NPV at the hurdle rate is nonnegative.
  \item Carried interest is $k$ calls struck at the hurdle less $k-c$ calls struck at the catch-up
        threshold. In the option model its sensitivity to risk is positive near the hurdle and can be
        negative far above it.
  \item Deal-by-deal carry is at least whole-fund carry without hurdles. The modelled deal-by-deal
        arrangement produces frequent clawback liabilities, and repayment at termination does not
        undo the cost of earlier carry payments.
  \item Hedge fund incentive fees have the payoff of options on NAV struck at the high-water mark. In the
        models here fees take a larger share of expected gains from riskier funds, and in the deterministic perpetuity model a fee
        stream discounted at the fund's gross return is worth the opening NAV; real valuations depend
        on survival, performance, expenses and the discount rate.
  \item A fund of funds pays incentive fees to winning managers without netting the losers, so the
        incentive charge on a positive portfolio gain can exceed the headline rate; whether it does
        depends on how the underlying returns are spread.
  \item Each pressure in the fee terms has a contractual answer, from the GP's own capital and the
        clawback to key person clauses and, as a last resort, removal of the GP.
  \item Promote hurdles based on IRR reward speed; hurdles based on multiples reward size.
  \item Commitment programmes reach a steady state equal to the sum of one fund's profile over its ages,
        which helps estimate the annual commitments and liquidity reserves needed to sustain a target
        allocation.
\end{itemize}

\FloatBarrier
% =====================================================================================
\section{Exercises}
\label{sec2:exercises}

\begin{enumerate}[label=\textbf{\thechapter.\arabic*}, leftmargin=*]
  \item\label{exr2:1} A fund with commitments of 200 charges 2\% of commitments for five years and then
        1.5\% of invested cost, which averages 120 over the following five years. Find total
        management fees and the amount of commitment left for investment.
  \item\label{exr2:2} (a) Show from \cref{eq2:catchup} that with a full catch-up and a single
        distribution the catch-up tier has size $y^* = c(H-K)/(1-c)$ and the GP's share of profit is
        exactly $c$ once $V \ge V^*$. (b) Using \cref{prop2:single}, find $V^*$ for $k = 0.5$ and $c = 0.2$ in
        terms of $H$ and $K$.
  \item\label{exr2:3} A fund calls 100 at time 0 and distributes 150 at year 4, with an 8\% hurdle
        and 20\% carry. Find the GP's carry with a full catch-up and with a hard hurdle.
  \item\label{exr2:4} A hedge fund starts at 100 with 1.5\% and 20\% fees and a high-water mark.
        Gross returns are 20\%, $-15\%$ and 12\%. Find the incentive fee and closing NAV each year.
  \item\label{exr2:5} Prove that $(a+b)^{+} \le a^{+} + b^{+}$ for all real $a$, $b$, and use it to
        explain why a fund of funds investor can pay incentive fees in a year when its portfolio loses
        money.
  \item\label{exr2:6} A fund of 500 earns 8\% gross with 1.5\% and 15\% fees and the manager discounts
        at 8\%. Use \cref{eq2:annuity} to value five years of fees.
  \item\label{exr2:7} Use \cref{eq2:vegapeak} to find the NAV at which the vega of a one-year
        incentive fee is greatest when $\sigma = 40\%$ and $r = 3\%$, as a percentage of the high-water
        mark. Explain what this means for a manager whose fund has fallen 15\%.
  \item\label{exr2:8} In \cref{ex2:ta}, verify the steady uncalled commitment of
        \progUncalledSteady{} from the closed form for $U_t$, and explain why it does not depend on
        $G$, $B$ or $L$.
\end{enumerate}

\subsection*{Computing projects}
\label{sec2:projects}
These use the data and code in the repository. Each is solved by a function in
\texttt{code/ch02/projects.py}, which \texttt{test\_projects.py} runs, so an answer can be
checked with \texttt{python code/ch02/projects.py} from the root of the repository. The
shipped series end in September 2026, and those that end earlier are noted.

\begin{enumerate}[label=\textbf{P2.\arabic*}, leftmargin=*, ref=P2.\arabic*]
\item Use \texttt{takahashi\_alexander} to trace a fund's cumulative net cash flow for annual growth of 6\%, 12\%
and 18\%. In which year is the J-curve lowest and in which does it turn positive? Explain why the trough
barely moves with growth while the break-even year can.
\item With \texttt{fee\_annuity\_pv}, compute the present value of ten years of 2 and 20 fees as a share of
the initial assets, for gross returns from 4\% to 12\% discounted at the gross return. Explain why the share
rises with the gross return although the discount rate rises with it.
\end{enumerate}

\subsection*{Answers}
\begin{enumerate}[label=\textbf{\thechapter.\arabic*}, leftmargin=*]
  \item $0.02\times200\times5 + 0.015\times120\times5 = 20 + 9 = \exOneA$; if all commitments are
        eventually called, $200 - \exOneA = \exOneB$ can be invested.
  \item (a) With $k=1$, $P=K$, $D=H$ and $G=0$, \cref{eq2:catchup} gives $y^* = c(H-K)/(1-c)$, so
        $V^* = H + y^*$. There the GP holds $y^* = c(H-K) + c\,y^* = c(V^*-K)$, and each later dollar
        adds $c$ to carry and one to profit, so the share stays $c$. (b) $V^* = (kH - cK)/(k-c) =
        (0.5H - 0.2K)/0.3 = (5H - 2K)/3$.
  \item $H = 100\times1.08^4 = \exThreeH$ and $V^* = (\exThreeH - 20)/0.8 = \exThreeVstar < 150$, so the
        GP receives $0.2\times50 = \exThreeGP$. With a hard hurdle it receives
        $0.2\times(150 - \exThreeH) = \exThreeGPHard$.
  \item Year 1: $120 - 1.5 = 118.5$, fee $0.2\times18.5 = \exFourIncOne$, NAV \exFourCloseOne{} (the new
        mark). Year 2: $\exFourCloseOne\times0.85 - 0.015\times\exFourCloseOne = \exFourCloseTwo$, no
        incentive fee. Year 3: $\exFourCloseTwo\times1.12 - 0.015\times\exFourCloseTwo = \exFourCloseThree$,
        below the mark, so no incentive fee.
  \item $a + b \le a^{+} + b^{+}$ and $0 \le a^{+} + b^{+}$, so $\max\{a+b, 0\} \le a^{+} + b^{+}$. If one
        underlying fund gains and another loses more, the portfolio loses, but the winning fund still
        charges $p$ times its own gain: the left side is zero and the right side is positive.
  \item $\phi = 0.015 + 0.15\times0.065 = 2.475\%$, $G = 0.85\times0.065 = 5.525\%$,
        $q = 1.05525/1.08$, and $\mathrm{PV}_5 = \frac{0.02475\times500}{1.08}\cdot\frac{1-q^5}{1-q} = \exSixPV$.
  \item $S^*/M = \exp(0.08 - 0.03) = \exSevenRatio$, about 105\% of the mark. At the peak the fee gains
        \exSevenVegaPeak{} per point of volatility per 100 of mark; at 85\% of the mark it still gains
        \exSevenVegaDown. At high volatility the vega curve is wide (panel (b) of
        \cref{fig2:optionview}), so a manager 15\% below the mark still has a strong incentive to add risk.
  \item $\sum_{a\ge0}U_a = 1 + 0.75 + 0.5 + 0.25\,(1 + 0.5 + 0.25 + \dots) = 2.25 + 0.5 = 2.75$. The
        uncalled commitment depends only on the contribution rates in \cref{eq2:taC}; the
        growth and distribution parameters do not enter it.
\end{enumerate}

\FloatBarrier
% =====================================================================================
\section{Notes, sources and credits}
\label{sec2:credits}

\paragraph{Sources}
The promote tier rates are illustrative assumptions with no market data behind them. The cash flow
modelling of \cref{sec2:pacing} anticipates the simulated funds of Chapter~11. The option
valuation uses the formula of \textcite{black1973}, derived in Chapter~5; \textcite{hull2022} and
\textcite{wilmott2006} give background. Statements about common terms, such as fund life, GP commitment, removal votes and hurdle rates,
describe general market practice; the empirical findings are cited where they appear. Further reading: \textcite{metrick2010} on the economics of private equity fees, \textcite{phalippou2009}
on fees and costs, \textcite{goetzmann2003} on high-water marks, \textcite{carpenter2000} on option pay and
risk, \textcite{agarwal2009} and \textcite{brown2001} on hedge fund incentives, \textcite{brown2003} on
fees in funds of funds, \textcite{takahashi2002} on commitment pacing and \textcite{ilpa2019} on investor
guidance for fund terms.

\paragraph{Open source projects}
The code uses NumPy \parencite{numpy}, SciPy \parencite{scipy} and Matplotlib \parencite{matplotlib}; pyxirr \parencite{pyxirr} provides an independent IRR
implementation used in the tests, the R cross-checks run on R \parencite{rcore} and the Octave cross-checks on GNU Octave
\parencite{octave}. The
full list, with links, is in \texttt{CREDITS.md}.

\end{chapterrefs}
